% VALEURS DE FONCTIONS L ET PERIODES D'INTEGRALES
% P. Deligne
% Proceedings of Symposia in Pure Mathematics, Vol. 33 (1979), part 2, pp. 313-346

\documentclass[12pt,a4paper]{article}
\usepackage{fontspec}
\usepackage[french]{babel}
\usepackage{amsmath, amssymb, amsthm, amscd}
\usepackage{mathtools}
\usepackage{enumitem}
\usepackage{graphicx}
\usepackage{xcolor}
\usepackage[margin=2.5cm]{geometry}
\usepackage{hyperref}
\usepackage{tikz-cd}
\usepackage{booktabs}
\usepackage{array}
\usepackage{stmaryrd}

\theoremstyle{plain}
\newtheorem{theorem}{Theoreme}[section]
\newtheorem{lemma}[theorem]{Lemme}
\newtheorem{proposition}[theorem]{Proposition}
\newtheorem{corollary}[theorem]{Corollaire}

\theoremstyle{definition}
\newtheorem{definition}{Definition}[section]
\newtheorem{remark}[definition]{Remarque}
\newtheorem{example}[definition]{Exemple}

\theoremstyle{remark}
\newtheorem{hope}{Espoir}
\newtheorem{conjecture}[definition]{Conjecture}
\newtheorem{proposition-definition}[theorem]{Proposition-Definition}

\DeclareMathOperator{\Hom}{Hom}
\DeclareMathOperator{\End}{End}
\DeclareMathOperator{\Aut}{Aut}
\DeclareMathOperator{\Gal}{Gal}
\DeclareMathOperator{\Spec}{Spec}
\DeclareMathOperator{\Sym}{Sym}
\DeclareMathOperator{\Ker}{Ker}
\DeclareMathOperator{\Image}{Im}
\DeclareMathOperator{\rank}{rang}
\DeclareMathOperator{\Res}{Res}
\DeclareMathOperator{\Tr}{Tr}
\DeclareMathOperator{\Ind}{Ind}
\DeclareMathOperator{\GL}{GL}
\DeclareMathOperator{\SL}{SL}
\DeclareMathOperator{\detred}{det^{red}}

\newcommand{\Q}{\mathbb{Q}}
\newcommand{\R}{\mathbb{R}}
\newcommand{\C}{\mathbb{C}}
\newcommand{\Z}{\mathbb{Z}}
\newcommand{\HH}{\mathbb{H}}
\newcommand{\A}{\mathbb{A}}
\newcommand{\PP}{\mathbb{P}}
\newcommand{\Fb}{\mathrm{F}_\infty}
\newcommand{\dpart}{\partial}

\begin{document}

\thispagestyle{empty}

\begin{center}
{\large\bfseries VALEURS DE FONCTIONS $L$ ET PERIODES D'INTEGRALES}\\[1em]
P. DELIGNE\\[1em]
Proceedings of Symposia in Pure Mathematics\\
Vol.~33 (1979), part~2, pp.~313--346
\end{center}

\vspace{1em}


Dans cet article, j'enonce une conjecture (1.8, 2.8) reliant les valeurs de
certaines fonctions $L$ en certains points entiers a des periodes d'integrales.

Les fonctions $L$ considerees sont celles des motifs---un mot auquel on n'attachera
pas un sens precis. Ceci inclut notamment les fonctions $L$ d'Artin, les fonctions $L$
attachees a des caracteres de Hecke algebriques (= Grossencharakter de type $A_0$),
et celles attachees aux formes modulaires holomorphes sur le demi-plan de Poincare,
supposees primitives (= new forms; on considere toutes les fonctions $L$,
$L_{Sym^n}$, attachees aux puissances symetriques, $Sym^n$, de la representation
$l$-adique correspondante).

Cet article doit le jour a D. Zagier: pour son insistance a demander une conjecture,
et pour la confirmation experimentale qu'il en a donnee, sitot emise, pour les
fonctions $L_3$ et $L_5$ attachees a $\Delta = \sum_{n=1}^{\infty} \tau(n)q^n$
(voir [18]). C'est cette confirmation qui m'a donne la confiance necessaire pour
verifier que la conjecture etait compatible aux resultats de Shimura [13] sur les
valeurs de fonctions $L$ de caracteres de Hecke algebriques.

\section*{0. Motifs}
\addcontentsline{toc}{section}{0. Motifs}

Le lecteur est invite a ne consulter ce paragraphe qu'au fur et a mesure des besoins.
On y rappelle une partie du formalisme, defini par Grothendieck, des motifs. Pour les
demonstrations, je renvoie a [8].

\subsection*{0.1}
\addcontentsline{toc}{subsection}{0.1}

La definition de Grothendieck des motifs sur un corps $k$ a la forme suivante.

Soit $\mathcal{V}(k)$ la categorie des varietes projectives et lisses sur $k$. On construit

(a) une categorie additive $\mathcal{M}'(k)$, pour laquelle les groupes $\Hom(M, N)$ sont des
espaces vectoriels sur $\Q$, munie de

\quad (i) Un produit tensoriel $\otimes$, associatif, commutatif et distributif par rapport a
l'addition des objets (``associatif'' et ``commutatif'' n'est pas une propriete du foncteur
$\otimes$, mais une donnee, soumise a certaines compatibilites---cf.~Saavedra [19]);

\quad (ii) Un foncteur contravariant $H^*$, de $\mathcal{V}(k)$ dans $\mathcal{M}'(k)$, bijectif sur les objets et
transformant sommes disjointes en sommes, et produits en produits tensoriels
(donnee d'un isomorphisme de foncteurs $H^*(X \times Y) = H^*(X) \otimes H^*(Y)$,
compatible a l'associativite et a la commutativite).

Il s'agit, pour l'essentiel, de definir $\Hom(H^*(X), H^*(Y))$. Pour la definition, utilisee
par Grothendieck, de ce groupe comme un groupe de classes de correspondances
entre $X$ et $Y$, voir 0.6.

(b) Rappelons (SGA 4 IV 7.5) qu'une categorie additive est karoubienne si tout
projecteur (= endomorphisme idempotent) est defini par une decomposition en
somme directe, et que chaque categorie additive a une enveloppe karoubienne,
obtenue en lui adjoignant formellement les images des projecteurs. La categorie
$\mathcal{M}^{eff}(k)$ des motifs effectifs sur $k$ est l'enveloppe karoubienne de $\mathcal{M}'(k)$.

(c) On definit le motif de Tate $\Z(-1)$ comme un facteur direct $H^2(\PP^1)$ convenable
de $H^*(\PP^1)$. On verifie que la symetrie (deduite de la commutativite de $\otimes$):
$\Z(-1) \otimes \Z(-1) \to \Z(-1) \otimes \Z(-1)$ est l'identite, et que le foncteur
$M \mapsto M \otimes \Z(-1)$ est pleinement fidele.

(d) La categorie $\mathcal{M}(k)$ des motifs sur $k$ se deduit de $\mathcal{M}^{eff}(k)$ en rendant inversible
le foncteur $M \mapsto M \otimes \Z(-1)$.

Notons $(n)$ l'itere $(-n)$-ieme de l'auto-equivalence $M \mapsto M \otimes \Z(-1)$. La categorie
$\mathcal{M}(k)$ admet $\mathcal{M}^{eff}(k)$ comme sous-categorie pleine, et tout objet de $\mathcal{M}(k)$ est
de la forme $M(n)$ pour $M$ dans $\mathcal{M}^{eff}(k)$ et $n$ un entier.

Par definition, si $F$ est un foncteur additif de $\mathcal{M}^{eff}(k)$ dans une categorie karoubienne
$\mathcal{A}$, il se prolonge a $\mathcal{M}^{eff}(k)$. Si $\mathcal{A}$ est munie d'une auto-equivalence $A \mapsto A(-1)$,
et le prolongement de $F$ d'un isomorphisme de foncteurs $F(M(-1)) = F(M)(-1)$,
il se prolonge a $\mathcal{M}(k)$.

\begin{example}
Si $k'$ est une extension de $k$, et $I^*$ le foncteur $H^*(X) \mapsto H^*(X \otimes_k k')$
de $\mathcal{M}^{eff}(k)$ dans $\mathcal{M}^{eff}(k')$, on obtient le foncteur extension des scalaires de
$\mathcal{M}(k)$ dans $\mathcal{M}(k')$. Si $k'$ est une extension finie de $k$, on dispose du foncteur de
restriction des scalaires a la Grothendieck
\[
\Res_{k'/k}: \mathcal{V}(k') \to \mathcal{V}(k): X \mapsto \Spec(k') \mapsto (X \to \Spec(k') \to \Spec(k)).
\]
On le prolonge en $F: H^*(X) \mapsto H^*(\Res_{k'/k} X)$, d'ou un foncteur de restriction
des scalaires $R_{k'/k}: \mathcal{M}(k') \to \mathcal{M}(k)$.
\end{example}

\begin{example}
Soit $\mathcal{H}$ une ``theorie de cohomologie'', a valeurs dans une categorie karoubienne
$\mathcal{A}$, fonctorielle pour les morphismes dans $\mathcal{M}^{eff}(k)$. Le foncteur $\mathcal{H}$ se prolonge
a $\mathcal{M}^{eff}(k)$. Si $\mathcal{A}$ est munie d'un produit tensoriel, que $\mathcal{H}$ verifie une formule de
Kunneth et que la tensorisation avec $\mathcal{H}(\Z(-1))$ est une auto-equivalence de $\mathcal{A}$,
il se prolonge a $\mathcal{M}(k)$. Ce prolongement est le foncteur ``realisation d'un motif
dans la theorie $\mathcal{H}$''.

Nous noterons $(n)$ l'auto-equivalence de $\mathcal{A}$ iteree $(-n)$-ieme de la tensorisation
avec $\mathcal{H}(\Z(-1))$. Pour la determination de $\mathcal{H}(\Z(-1))$ dans diverses theories $\mathcal{H}$,
voir 3.1.
\end{example}

\subsection*{0.2}
\addcontentsline{toc}{subsection}{0.2}

Nous utiliserons les realisations suivantes:

\subsubsection*{0.2.1. Realisation de Betti $H_B$}
\addcontentsline{toc}{subsubsection}{0.2.1}

Correspondant a $k \subset \C$, $\mathcal{A} =$ espaces vectoriels sur $\Q$,
$\mathcal{H} =$ cohomologie rationnelle: $X \mapsto H^*(X(\C), \Q)$;

\subsubsection*{0.2.2. Realisation de de Rham $H_{DR}$}
\addcontentsline{toc}{subsubsection}{0.2.2}

Correspondant a $k$ de caracteristique $0$, $\mathcal{A} =$ espaces vectoriels sur $k$,
$\mathcal{H} =$ cohomologie de de Rham: $X \mapsto H^*(X, \Omega^*_X)$;

\subsubsection*{0.2.3. Realisation $l$-adique $H_l$}
\addcontentsline{toc}{subsubsection}{0.2.3}

Correspondant a $k$ algebriquement clos, de caracteristique $\neq l$,
$\mathcal{A} =$ espaces vectoriels sur $\Q_l$, $\mathcal{H} =$ cohomologie $l$-adique:
$X \mapsto H^*(X, \Q_l)$;

Et leurs variantes:

\subsubsection*{0.2.4. Realisation de Hodge}
\addcontentsline{toc}{subsubsection}{0.2.4}

$k$ est ici une cloture algebrique de $\R$, et $\mathcal{A}$ la categorie des espaces vectoriels
sur $\Q$, de complexifie $V \otimes_{\Q} k$ muni d'une bigraduation
$V \otimes_{\Q} k = \bigoplus V^{p,q}$ telle que $V^{q,p}$ soit le complexe conjugue de $V^{p,q}$.
Pour theorie de cohomologie, on prend le foncteur $X \mapsto H^*(X(k), \Q)$ muni de la
bigraduation de son complexifie $H^*(X(k), k)$ fournie par la theorie de Hodge. La
realisation de Betti est sous-jacente a celle de Hodge.

\subsubsection*{0.2.5}
\addcontentsline{toc}{subsubsection}{0.2.5}

Pour $k$ quelconque, et $\sigma$ un plongement complexe de $k$, on note $H_\sigma(M)$ la
realisation de Betti du motif sur $\C$ deduit de $M$ par l'extension des scalaires
$\sigma: k \to \C$. Notons $c$ la conjugaison complexe. On obtient par transport de structure
un isomorphisme $F_\sigma: H_\sigma(M) \xrightarrow{\sim} H_{c\sigma}(M)$, et $F_\sigma \otimes c$ envoie
$H_\sigma^{p,q}$ sur $H_{c\sigma}^{q,p}$. Pour $\sigma$ reel, $F_\sigma$ est une involution de $H_\sigma(M)$, dont le
complexifie echange $H_\sigma^{p,q}$ et $H_\sigma^{q,p}$.

La $\sigma$-realisation de Hodge est $H_\sigma(M)$, muni de sa bigraduation de Hodge et,
si $\sigma$ est reel, de l'involution $F_\sigma$. Pour $k = \Q$, $k \subset \R$ et le plongement
identique, on remplace l'indice $\sigma$ par $B$. Pour $k = \R$, et $M = H^*(X)$, l'involution
$F_\sigma$ de $H_B(M) = H^*(X(\C), \Q)$ est l'involution induite par la conjugaison complexe
$F_\infty: X(\C) \to X(\C)$.

\subsubsection*{0.2.6. Realisation de de Rham}
\addcontentsline{toc}{subsubsection}{0.2.6}

La cohomologie de de Rham est munie d'une filtration naturelle, la filtration de Hodge,
aboutissement de la suite spectrale d'hypercohomologie
\[
E_1^{pq} = H^q(X, \Omega_X^p) \Rightarrow H_{DR}^{p+q}(X).
\]
D'ou, un filtration $F$ sur $H_{DR}(M)$.

\subsubsection*{0.2.7. Realisation $l$-adique}
\addcontentsline{toc}{subsubsection}{0.2.7}

Si $X$ est une variete sur un corps $k$, de cloture algebrique $\overline{k}$, le groupe de Galois
$\Gal(\overline{k}/k)$ agit par transport de structure sur $H^*(X_{\overline{k}}, \Q_l)$. Si $M$ est un motif
sur $k$, et qu'on definit $H_l(M)$ comme la realisation $l$-adique du motif sur $\overline{k}$ qui
s'en deduit par extension des scalaires, ceci fournit une action de $\Gal(\overline{k}/k)$ sur $H_l(M)$.

\subsubsection*{0.2.8. Realisation adelique finie}
\addcontentsline{toc}{subsubsection}{0.2.8}

Pour $k$ algebriquement clos, de caracteristique $0$, on peut rassembler les cohomologies
$l$-adiques en une cohomologie adelique
\[
H^*(X, \A^f) = (\prod_l H^*(X, \Z_l)) \otimes_{\Z} \Q.
\]
On peut cumuler les variantes 0.2.7 et 0.2.8.

\subsubsection*{0.2.9}
\addcontentsline{toc}{subsubsection}{0.2.9}

Dans les exemples 0.2.1, 0.2.2, 0.2.3, et leurs variantes, on peut remplacer la categorie
$\mathcal{A}$ par celle, $\mathcal{A}^{gr}$, des objets gradues de $\mathcal{A}$: utiliser la graduation naturelle de $H^*$.
Pour que la formule de Kunneth fournisse un isomorphisme de foncteurs
$\mathcal{H}(M \otimes N) \xrightarrow{\sim} \mathcal{H}(M) \otimes \mathcal{H}(N)$, compatible a l'associativite et a la
commutativite, il faut prendre pour donnee de commutativite dans $\mathcal{A}^{gr}$ celle donnee par
la regle de Koszul.

\subsection*{0.3}
\addcontentsline{toc}{subsection}{0.3}

Pour $k \subset \C$, le theoreme de comparaison entre cohomologie classique et
cohomologie etale fournit un isomorphisme
$H^*(X(\C), \Q) \otimes_{\Q} \Q_l \xrightarrow{\sim} H^*(X, \Q_l)$.
Si $X$ est defini sur $\R$, cet isomorphisme transforme $F_\infty$ (0.2.5) en l'action (0.2.7) de
la conjugaison complexe. Pour un motif, on a de meme $H_B(M) \otimes \Q_l \xrightarrow{\sim} H_l(M)$.

\subsection*{0.4}
\addcontentsline{toc}{subsection}{0.4}

Pour $k \subset \C$, le complexe de de Rham holomorphe sur $X^{an}$ est une resolution du
faisceau constant $\C$. Par GAGA, on a donc un isomorphisme
\[
H^*(X(\C), \Q) \otimes \C = H^*(X(\C), \C) \xrightarrow{\sim} H^*(X^{an}, \Omega_{X^{an}}^*) \xleftarrow{\sim} H^*(X, \Omega_X^*).
\]
Pour un motif, on obtient un isomorphisme (compatible aux filtrations de Hodge:
$F^q = \bigoplus_{p' \geq q} V^{p',q'}$, cf.~0.2.4 et 0.2.6)
\begin{equation}\label{eq:0.4.1}
I: H_B(M) \otimes_{\Q} \C \xrightarrow{\sim} H_{DR}(M) \otimes_{k,\sigma} \C.
\end{equation}

Soit plus generalement $M$ un motif sur un corps $k$, et $\sigma$ un plongement complexe de $k$.
Appliquant ce qui precede au motif sur $\C$ deduit de $M$ par extension des scalaires, on
trouve un isomorphisme
\[
I: H_\sigma(M) \otimes_{\Q} \C \xrightarrow{\sim} H_{DR}(M) \otimes_{k,\sigma} \C.
\]

Notons le cas particulier $k = \Q$, $\sigma$ le plongement identique: on trouve deux
$\Q$-structures naturelles sur la realisation de $M$ en cohomologie complexe: l'une,
$H_B(M)$, liee a la description de la cohomologie en termes de cycles, l'autre,
$H_{DR}(M)$, liee a sa description en termes de formes differentielles algebriques.

\subsection*{0.5}
\addcontentsline{toc}{subsection}{0.5}

Ce qu'il advient de ces realisations et compatibilites par extension du corps des scalaires
est clair. Pour la restriction des scalaires, les resultats sont les suivants.

Soient $k'$ une extension finie de $k$, $M'$ un motif sur $k'$, et $M = R_{k'/k}(M')$
\begin{equation}\label{eq:0.5.1}
H_\sigma(M) = \bigoplus_{\sigma'} H_{\sigma'}(M'),
\end{equation}
o\`u la somme est etendue a l'ensemble $I(\sigma')$ des plongements complexes de $k'$ qui
prolongent $\sigma$. Cet isomorphisme est compatible aux bigraduations de Hodge, et a
$F_\infty$.
\begin{equation}\label{eq:0.5.2}
H_{DR}(M) = H_{DR}(M') \quad \text{(restriction des scalaires de } k' \text{ a } k \text{).}
\end{equation}
Cet isomorphisme est compatible a la filtration de Hodge.
\begin{equation}\label{eq:0.5.3}
H_l(M) = \Ind H_l(M') \quad \text{(representation induite de } \Gal(\overline{k}/k') \text{ a } \Gal(\overline{k}/k) \text{).}
\end{equation}

Via l'isomorphisme $k' \otimes_{k,\sigma} \C = \C^{I(\sigma')}$, et l'isomorphisme
$H_{DR}(M') \otimes_{k',\sigma'} \C = \bigoplus_{\sigma'} H_{DR}(M') \otimes_{k',\sigma'} \C$
qui s'en deduit, (0.4.1) pour $M$ et $\sigma$ est la somme des isomorphismes (0.4.1) pour
$M'$ et $\sigma'$ ($\sigma' \in I(\sigma')$):
\[
\begin{CD}
H_\sigma(M) \otimes_{\Q} \C @>{\text{(0.4.1)}}>> H_{DR}(M) \otimes_{k,\sigma} \C \\
@| @| \\
\bigoplus_{\sigma'} H_{\sigma'}(M') \otimes_{\Q} \C @>{\oplus \text{(0.4.1)}}>>
\bigoplus_{\sigma'} H_{DR}(M') \otimes_{k',\sigma'} \C
\end{CD}
\]


\subsection*{0.6}
\addcontentsline{toc}{subsection}{0.6}

Pour $X$ projectif et lisse sur $k$, notons $Z^d(X)$ l'espace vectoriel sur $\Q$ de base
l'ensemble des sous-schemas fermes irreductibles de codimension $d$ de $X$, et $Z^d_\sim(X)$
son quotient par une relation d'equivalence $R$. Pour $k$ de caracteristique $0$, une des
definitions de Grothendieck des motifs s'obtient en faisant (pour $X$ et $Y$ connexes)
$\Hom(H^*(Y), H^*(X)) = Z^{\dim Y}_\sim(X \times Y)$, $R$ etant l'equivalence cohomologique.
En ce qui concerne la caracteristique $p$, signalons seulement deux difficultes:
on ignore si l'equivalence cohomologique---en cohomologie $l$-adique ($l \neq p$)---
est independante de $l$, et la definition de la classe d'un cycle pose des problemes en
cohomologie cristalline.

\subsection*{0.7}
\addcontentsline{toc}{subsection}{0.7}

Soit $X$ une variete projective et lisse complexe. Nous appellerons \emph{cycle de Hodge}
de codimension $d$ sur $X$ un element de $H^{2d}(X(\C), \Q)$ de type $(d,d)$---soit, ce qui
revient au meme, un element de $H^{2d}(X(\C), \Q)(d)$ de type $(0,0)$.

Soient $k$ un corps algebriquement clos et $X$ une variete projective et lisse sur $k$.
Posons
\[
H^{2d}(X, k \times \A^f)(d) = H^{2d}_{et}(X, \A^f) \otimes_{\Q} H^{2d}(X, \A^f)(d).
\]
Pour $k = \C$, on appelle encore cycle de Hodge l'image dans
\[
H^{2d}(X, k \times \A^f)(d) = H^{2d}(X(\C), \Q)(d) \otimes_{\Q} (k \times \A^f)
\]
d'un cycle de Hodge. Pour $k$ algebriquement clos, admettant des plongements complexes,
un \emph{cycle de Hodge absolu} de codimension $d$ sur $X$ est un element de
$H^{2d}(X, k \times \A^f)(d)$, tel que, pour tout plongement complexe $\sigma$ de $k$, son image
dans $H^{2d}(X \otimes_k \C, \C \times \A^f)(d)$ soit de Hodge. On verifie que

\begin{proposition}\label{prop:0.8}
\textup{(i)} L'espace vectoriel sur $\Q$ $Z^d_{\text{Ha}}(X)$ des cycles de Hodge absolu est
invariant par extension des scalaires de $k$ a un corps algebriquement clos $k'$ (admettant
encore un plongement complexe).

\textup{(ii)} Pour $k$ algebriquement clos de caracteristique $0$, et $X$ defini sur un sous-corps
algebriquement clos $k_0$ de $k$, admettant un plongement complexe:
$X = X_0 \otimes_{k_0} k$, on pose
\[
Z^d_{\text{Ha}}(X_0) = Z^d_{\text{Ha}}(X) \subset H^{2d}(X, k \times \A^f)(d) \xrightarrow{\sim} H^{2d}(X_0, k \times \A^f)(d).
\]
D'apres (i), cette definition ne depend pas des choix de $X_0$ et $k_0$.

\textup{(iii)} Pour $X$ defini sur un sous-corps $k_0$ de $k$: $X = X_0 \otimes_{k_0} k$, le groupe
$\Aut(k/k_0)$ agissant sur $H^{2d}(X, k \times \A^f)(d)$ stabilise $Z^d_{\text{Ha}}(X)$. Il agit sur
$Z^d_{\text{Ha}}(X)$ a travers un groupe fini, correspondant a une extension finie $k_1$ de $k_0$.
On pose $Z^d_{\text{Ha}}(X_0) = Z^d_{\text{Ha}}(X)^{\Aut(k/k_0)}$.
\end{proposition}

\subsection*{0.9}
\addcontentsline{toc}{subsection}{0.9}

Une notion utile de motif s'obtient en faisant (pour $X$ et $Y$ connexes)
\[
\Hom(H^*(Y), H^*(X)) = Z^{\dim Y}_{\text{Ha}}(X \times Y).
\]
Les composantes de Kunneth de la diagonale de $X \times X$ sont absolument de Hodge.
Ceci permet de decomposer $H^*(X)$ en une somme de motifs $H^i(X)$, de munir la
categorie des motifs d'une graduation (avec $H^i(X)$ de poids $i$) et de modifier la
contrainte de commutativite pour $\otimes$ comme en [19, VI, 4.2.1.4]. On verifie que, ceci
fait, la $\otimes$-categorie des motifs sur $k$ est tannakienne, isomorphe a la categorie des
representations d'un groupe proalgebrique reductif.

Pour $k = \C$, la conjecture suivante, plus faible que celle de Hodge, equivaut a dire
que le foncteur ``realisation de Hodge'' est une equivalence de la categorie des motifs
0.9 avec la categorie des structures de Hodge facteur direct de la cohomologie de
varietes algebriques (ou deduites par twist a la Tate de tels facteurs directs).

\begin{hope}\label{esp:0.10}
Tout cycle de Hodge l'est absolument.
\end{hope}

Si $X$ est une variete abelienne a multiplication complexe par un corps quadratique
imaginaire $K$, avec $\text{Lie}(X)$ libre sur $k \otimes K$, les methodes de B.~Gross [7]
prouvent que certains cycles de Hodge non triviaux sont absolument de Hodge.

Ajoute sur epreuves: partant de ce cas particulier, j'ai pu verifier 0.10 pour les varietes
abeliennes. J'ai verifie aussi que la $\otimes$-categorie de motifs (0.9) engendree par les
$H^n(X)$, pour $X$ une variete abelienne, contient les motifs $H^*(X)$, pour $X$ une surface
K3 ou une hypersurface de Fermat.

\subsection*{0.11}
\addcontentsline{toc}{subsection}{0.11}

Le principal defaut de la definition 0.9 des motifs est qu'elle ne se prete pas a la
reduction mod $p$. On ignore si un motif 0.9 sur un corps de nombres $F$ fournit un
systeme compatible de representations $l$-adiques de $\Gal(\overline{F}/F)$.

\subsection*{0.12}
\addcontentsline{toc}{subsection}{0.12}

Nous utiliserons le mot ``motif'' de facon libre, sans nous preoccuper de faire rentrer
les motifs consideres dans le cadre de Grothendieck. L'essentiel pour nous sera de
disposer de realisations $\mathcal{H}(M)$, pour les theories $\mathcal{H}$ considerees en 0.2, et d'avoir pour
ce groupes le meme formalisme que pour les $H^*(X)$.


\section*{1. Enonce de la conjecture (cas rationnel)}
\addcontentsline{toc}{section}{1. Enonce de la conjecture (cas rationnel)}

\subsection*{1.1}
\addcontentsline{toc}{subsection}{1.1}

Soit $M$ un motif sur $\Q$. Nous admettrons que les realisations $l$-adiques $H_l(M)$ de $M$
forment un systeme strictement compatible de representations $l$-adiques, au sens de
Serre [11, 1.1]. A savoir: il existe un ensemble fini $S$ de nombres premiers, tel que
chaque $H_l(M)$ soit non ramifie en dehors de $S \cup \{l\}$, et que, notant par
$F_p \in \Gal(\overline{\Q}/\Q)$ un element de Frobenius geometrique en $p$ (l'inverse d'une
substitution de Frobenius usuelle), le polynome $\det(1 - F_p t, H_l(M)) \in \Q_l[t]$
($p \notin S \cup \{l\}$) soit a coefficients rationnels, et independant de $l$. Notons
$Z_p(M, t)$ son inverse, et posons $L_p(M, s) = Z_p(M, p^{-s})$.

La serie de Dirichlet a coefficients rationnels donnee par le produit eulerien
\[
L_S(M, s) = \prod_{p \notin S} L_p(M, s)
\]
converge pour $\Re s$ assez grand. Pour $s$ quelconque, on definit $L_S(M, s)$ par un
prolongement analytique (qu'on espere exister).

Notre but est d'enoncer une conjecture donnant la valeur de $L_S(M, s)$ en certains
points entiers, au produit pres par un nombre rationnel. Puisque, pour $s$ entier,
$p^{-s}$ est rationnel, le choix de $S$ est sans importance---excepte qu'agrandir $S$ peut
introduire des zeros malvenus.

\subsection*{1.2}
\addcontentsline{toc}{subsection}{1.2}

Pour ecrire proprement l'equation fonctionnelle conjecturale des fonctions $L$, il y a
lieu de completer le produit eulerien $L_S(M, s)$ par des facteurs locaux en $p \in S$,
et a l'infini. La definition des facteurs locaux en $p \in S$ requiert une hypothese
additionnelle, que nous supposerons verifiee:

\begin{equation}\label{eq:1.2.1}
\text{Soient } p \text{ un nombre premier, } D_p \subset \Gal(\overline{\Q}/\Q) \text{ un groupe de decomposition}
\text{ en } p,\ I_p \subset D_p \text{ son sous-groupe d'inertie et } F_p \in D_p \text{ un Frobenius}
\text{ geometrique. Le polynome } \det(1 - F_p t, H_l(M)^{I_p}) \in \Q_l[t]\ (l \neq p)
\text{ est a coefficients rationnels, et est independant de } l.
\end{equation}

Posons $Z_p(M, t) = \det(1 - F_p t, H_l(M)^{I_p})^{-1} \in \Q(t)$, et $L_p(M, s) = Z_p(M, p^{-s})$.
On definit
\begin{equation}\label{eq:1.2.2}
L(M, s) = \prod_p L_p(M, s).
\end{equation}

Le facteur a l'infini $L_\infty(M, s)$ (essentiellement un produit de fonctions $\Gamma$)
depend de la realisation de Hodge de $M$---en fait seulement de la classe d'isomorphie de
l'espace vectoriel complexe $H_B(M) \otimes \C$, muni de sa decomposition de Hodge et
de l'involution $F_\infty$. Sa definition est rappelee en 5.2.

Posant $\Lambda(M, s) = L_\infty(M, s) L(M, s)$, l'equation fonctionnelle conjecturale des
fonctions $L$ s'ecrit
\begin{equation}\label{eq:1.2.3}
\Lambda(M, s) = \varepsilon(M, s) \cdot L(\check{M}, 1-s)
\end{equation}
o\`u $\check{M}$ est le dual de $M$ (de realisations les duales des realisations de $M$) et o\`u
$\varepsilon(M, s)$, comme fonction de $s$, est le produit d'une constante par une exponentielle.
Sa definition est rappelee en 5.2. Elle depend d'une hypothese additionnelle.

\begin{definition}\label{def:1.3}
Un entier $n$ est \emph{critique} pour $M$ si $L_\infty(M, s)$ et $L_\infty(\check{M}, 1-s)$ n'ont pas
de pole en $s = n$.
\end{definition}

Notre but est de conjecturer la valeur de $L(M, n)$ pour $n$ critique, a multiplication par
nombre rationnel pres. On a $L(M(n), s) = L(M, n+s)$ (3.1.2), et de meme pour $L_\infty$.
Ceci nous permet de ne considerer que les nombres $L(M) =_{dfn} L(M, 0)$. Nous dirons que
$M$ est critique si $0$ est critique pour $M$. On verifie que pour que $M$ soit critique, il faut
et il suffit que les nombres de Hodge $h^{p,q} =_{dfn} \dim_\C H_\sigma^{p,q}(M)$ de $M$, pour
$p \neq q$, ne soient non nuls que pour $(p,q)$ dans la partie hachuree du diagramme
ci-dessous, et que $F_\infty$ agisse sur $H_B^{p,p}(M)$ par l'identite si $p < 0$, par $-1$ si
$p \geq 0$:
\begin{equation}\label{eq:1.3.1}
\begin{matrix}
\\
\end{matrix}
\qquad
\begin{tikzpicture}[scale=0.5,baseline]
\fill[gray!30] (-4,4) -- (0,0) -- (4,0) -- (4,4) -- cycle;
\fill[gray!30] (-4,-4) -- (0,0) -- (-4,0) -- cycle;
\draw[->] (-4.5,0) -- (4.5,0) node[right] {$p$};
\draw[->] (0,-4.5) -- (0,4.5) node[above] {$q$};
\draw (-4,4) -- (4,-4);
\end{tikzpicture}
\end{equation}

Supposons $M$ homogene de poids $w$: $h^{p,q} = 0$ pour $p + q \neq w$, et posons
$\mathfrak{F}(M) = -w/2$. Il resulte de la conjecture de Weil que, pour $S$ assez grand, la serie
de Dirichlet $L_S(M, s)$ converge absolument pour $\mathfrak{F}(M) + \Re(s) > 1$. Pour
$\mathfrak{F}(M) + \Re(s) = 1$, on est au bord du demi-plan de convergence, et on conjecture

(a) que $L_S(M, s)$ ne s'annule pas,

(b) que $L_S(M, s)$ est holomorphe, sauf si $M$ est de poids pair $-2n$, et contient
$\Z(n)$ en facteur: on s'attend alors a un pole en $s = 1-n$ (une valeur non critique).

L'analogie avec le cas des corps de fonctions, et les cas connus, menent a croire
que les facteurs locaux $L_p(M, s)$ ($p$ quelconque, y compris $\infty$) n'ont de pole que
pour $\mathfrak{F}(M) + \Re(s) \leq 0$. Si tel est le cas, (a), (b) et l'equation fonctionnelle
conjecturale (1.2.3) impliquent que $L(M) \neq 0, \infty$ pour $M$ critique et $\mathfrak{F}(M) \neq \tfrac{1}{2}$.
Pour $\mathfrak{F}(M) = \tfrac{1}{2}$, $L(M)$ s'annule parfois. Notre conjecture (1.8) est alors vide.

\begin{proposition}\label{prop:1.4}
Soit $M$ un motif sur $\R$. Via l'isomorphisme (0.4.1)
\[
H_B(M) \otimes_\Q \C \xrightarrow{\sim} H_{DR}(M) \otimes_\R \C
\]
$H_{DR}(M)$ s'identifie au sous-espace de $H_B(M) \otimes \C$ fixe par $c \mapsto F_\infty \bar{c}$.
\end{proposition}

Prenons $M = H^*(X)$. La conjugaison complexe sur $H_{DR}(M)_\C = H^*(X_\C, \Omega^*_{X_\C})$
se deduit par fonctorialite de l'automorphisme antilineaire $F_\infty$ du schema $X_\C$.
Remontant la fleche composee (0.4) qui definit $I$, on l'identifie a l'involution deduite
de l'automorphisme $(F_\infty, F_\infty)$ de $(X(\C), \mathcal{O}^{an})$, puis au compose de
$F_\infty^*: H^*(X(\C), \C) \to H^*(X(\C), \C)$ et de la conjugaison complexe sur les coefficients.
Ceci verifie 1.4.

\subsection*{1.5}
\addcontentsline{toc}{subsection}{1.5}

Pour $M$ un motif sur $\R$, nous noterons $H^\pm_B(M)$ (resp.~$H^\pm_{DR}(M)$) le sous-espace
de $H_B(M)$ (resp.~$H_{DR}(M)$) fixe par $F_\infty$ (resp.~on $F_\infty = \pm 1$). Posons
$d^\pm(M) = \dim H^\pm_B(M)$ et $d^\pm_{DR}(M) = \dim H^\pm_{DR}(M)$. Pour $M$ un motif sur $k$,
et $\sigma$ un plongement complexe de $k$ qui se factorise par $\R$, on note $H^\pm_\sigma(M)$,
$d^\pm_\sigma(M)$ et $d^\pm_{DR}(M)$ ces objets pour le motif sur $\R$ deduit de $M$ par $\sigma$.
Pour $k = \Q$, on omet la mention de $\sigma$.

Le corollaire suivant resulte aussitot de 1.4, et de ce que, taut $F_\infty$ que la conjugaison
complexe echangent $H^{p,q}$ et $H^{q,p}$.

\begin{corollary}\label{cor:1.6}
Soit $M$ un motif sur $\R$. Pour $\mathfrak{F}$ la structure reelle $H_{DR}(M)$ de
$H_B(M) \otimes_\Q \C$, les sous-espaces $H^{p,q}_\mathfrak{F}$ sont definis sur $\R$; le sous-espace
$H^+_{DR}(M) \otimes_\R \C$ de $H_B(M) \otimes_\Q \C$ est reel, et $H^-_{DR}(M) \otimes_\R \C$ est
purement imaginaire.
\end{corollary}

\subsection*{1.7}
\addcontentsline{toc}{subsection}{1.7}

Dans la fin de ce paragraphe, nous ne considererons que des motifs sur $\Q$;
sauf mention expresse du contraire, nous les supposerons homogenes. Si leur poids
$w$ est pair, nous supposerons aussi que $F_\infty$ agit sur $H^{p,p}_B(M)$ ($w = 2p$) comme un
scalaire: soit $+1$, soit $-1$. Cette hypothese est verifiee pour $M$ critique. Puisque
$F_\infty$ echange $H^{p,q}_B$ et $H^{q,p}_B$, elle assure que les dimensions $d^+(M)$ et $d^-(M)$
sont egales l'une a $\sum_{p < q} h^{p,q}$, l'autre a $\sum_{p \geq q} h^{p,q}$. En particulier, ces dimensions
sont egales a celles de sous-espaces $F^+$ et $F^-$ figurant dans la filtration de Hodge.

Posant $H^\pm_{DR}(M) = H_{DR}(M)/F^\mp$, on a encore $\dim H^\pm_{DR}(M) = d^\pm(M)$.

Il resulte de ce que $F_\infty$ echange $H^{p,q}$ et $H^{q,p}$ que les applications composees
\begin{equation}\label{eq:1.7.1}
I^\pm: H^\pm_B(M) \hookrightarrow H_B(M) \xrightarrow{I} H_{DR}(M) \twoheadrightarrow H^\pm_{DR}(M)
\end{equation}
sont des isomorphismes. On pose
\begin{align}
c^\pm(M) &= \det(I^\pm), \label{eq:1.7.2}\\
\delta(M) &= \det(I). \label{eq:1.7.3}
\end{align}
Le determinant etant calcule dans des bases rationnelles de $H^\pm_B$ et $H^\pm_{DR}$
(resp.~$H_B$ et $H_{DR}$). La definition de $\delta(M)$ ne requiert pas les hypotheses faites sur $M$.

D'apres 1.6, $I^+$ est reel, i.e.~induit $I^+: H^+_B(M)_\R \xrightarrow{\sim} H^+_{DR}(M)_\R$,
tandis que $I^-$ est purement imaginaire. Les nombres $c^+(M)$, $i^{d^-(M)} c^-(M)$ et
$i^{d^-(M)} \delta(M)$ sont donc reels non nuls. A multiplication par un nombre rationnel
pres, ils ne dependent que de $M$.

Les periodes de $M$ sont classiquement les $(\omega, c)$ pour $\omega \in H_{DR}(M)$ et
$c \in H_B(M)^\vee$ (ou encore, si $X$ est une variete algebrique sur $\Q$, definies sur $\Q$
et que $c$ est un $n$-cycle sur $X(\C)$, $(\omega, c) = \int_c \omega$ est une periode de $H^n(X)$).
Exprimons $c^\pm(M)$ en terme de periodes. Le dual de $H^\pm_B(M)$ est le sous-espace
$F^\mp$ de $H_{DR}(M^\vee)$, ou $M^\vee$ est le motif dual de $M$. Si la base choisie de
$H^\pm_{DR}(M)$ a pour duale la base $(\omega_i)$ de $F^\mp(H_{DR}(M^\vee))$, et que $(c_i)$ est la
base choisie de $H^\pm_B(M)$, la matrice de $I^\pm$ est $(\omega_i, c_j)$, et
$c^\pm(M) = \det((\omega_i, c_j))$.

\begin{conjecture}\label{conj:1.8}
Si $M$ est critique, $L(M)$ est un multiple rationnel de $c^+(M)$.
\end{conjecture}


\section*{2. Enonce de la conjecture (cas general)}
\addcontentsline{toc}{section}{2. Enonce de la conjecture (cas general)}

\subsection*{2.1}
\addcontentsline{toc}{subsection}{2.1}

Jusqu'ici, nous n'avons considere que des fonctions $L$ donnees par des series de
Dirichlet a coefficients rationnels. Pour faire mieux, il nous faudra considerer des
motifs a coefficient dans des corps de nombres.

Voici deux facons, equivalentes, pour construire la categorie des motifs sur $k$,
a coefficient dans un corps de nombres $E$, a partir de la categorie des motifs sur $k$.
Il s'agit d'une construction valable pour toute categorie additive karoubienne
dans laquelle les $\Hom(X, Y)$ sont des espaces vectoriels sur $\Q$.

\textbf{A.} Un motif sur $k$, a coefficient dans $E$, est un motif $M$ sur $k$ muni d'une structure
de $E$-module: $E \to \End(M)$.

\textbf{B.} La categorie $\mathcal{M}(k)_E$ des motifs sur $k$, a coefficient dans $E$, est l'enveloppe
karoubienne (cf.~0.1(b)) de la categorie d'objets les motifs sur $k$---considere
comme objet de $\mathcal{M}^{gr}$, un motif $M$ se notera $M_E$---et de morphismes donnes par
$\Hom(X_E, Y_E) = \Hom(X, Y) \otimes_\Q E$.

\textit{Passage de B a A.} Pour $X$ un motif, et $V$ un espace vectoriel sur $\Q$ de dimension
finie, on note $X \otimes V$ le motif, isomorphe a une somme de $\dim(V)$ copies de $X$,
caracterise par $\Hom(Y, X \otimes V) = \Hom(Y, X) \otimes_\Q V$ (ou par
$\Hom(X \otimes V, Y) = \Hom(V, \Hom(X, Y))$). On passe de B a A en associant a
$M_E$ le motif $M \otimes E$, muni de sa structure de $E$-module naturelle.

\textit{Passage de A a B.} Si $M$ est muni d'une structure de $E$-module, on recupere sur
$M_E$ deux structures de $E$-module: celle deduite de celle de $M$, et celle qu'a tout
objet de $\mathcal{M}(k)_E$. L'objet de $\mathcal{M}(k)_E$ correspondant a $M$ est le plus grand facteur
direct de $M_E$ sur lequel ces structures coincident. En detail: l'algebre $E \otimes_\Q E$ est un
produit de corps, parmi lesquels une copie de $E$ dans laquelle $x \otimes 1$ et $1 \otimes x$ se
projettent tous deux comme $x$. L'idempotent correspondant agit sur $M_E$ (qui est
un $E \otimes_\Q E$-module) et son image est l'objet de $\mathcal{M}(k)_E$ qui correspond a $M$.

Les motifs a coefficient dans $E$ sont le plus souvent donnes sous la forme A. La
forme B a l'avantage de rester raisonnable pour $E$ non de rang fini sur $\Q$. Elle est
utile pour comprendre le formalisme tensoriel: on peut definir produit tensoriel et
dual, pour les motifs a coefficient dans $E$, par leur fonctorialite et les formules
$X_E \otimes_E Y_E = (X \otimes Y)_E$ et $(X_E)^\vee = (X^\vee)_E$. Dans le langage A, $X \otimes_E Y$
est le plus grand facteur direct de $X \otimes Y$ sur lequel coincident les deux structures de
$E$-module de $X \otimes Y$, et $X^\vee_E$ est le dual usuel de $X^\vee$, muni de la structure de
$E$-module transposee. Si on applique ces remarques a la categorie des espaces vectoriels
sur $\Q$, plutot qu'a celle des motifs, on retrouve l'isomorphisme du $E$-dual d'un espace
vectoriel sur $E$ avec son $\Q$-dual, donne par
\[
\omega \mapsto \text{la forme } \Tr_{E/\Q}(\omega, v).
\]

Nous avons defini en 0.1.1 des foncteurs de restriction et d'extension du corps $k$
des scalaires. Ils transforment motifs a coefficient dans $E$ en motifs a coefficient dans
$E$. On dispose aussi de foncteurs de restriction et d'extension des coefficients: soit
$F$ une extension finie de $E$:

\textit{Extension des coefficients.} Dans le langage A, c'est $X \mapsto X \otimes_E F$; dans le
langage B, c'est $X_E \mapsto X_F$.

\textit{Restriction des coefficients.} Dans le langage A, on restreint a $E$ la structure de
$F$-module.

Le lecteur prendra soin de ne pas confondre les roles de $k$ et $E$. Un exemple type
qu'on peut retenir est celui du $H^1$ des varietes abeliennes sur $k$, a multiplication
complexe par un ordre de $E$. En terme de la variete abelienne---prise a isogenie
pres---les foncteurs ci-dessus deviennent: extension du corps de base, restriction
des scalaires a la Weil, construction $\otimes_E F$, qui multiplie la dimension par $[F:E]$,
restriction a $E$ de la structure de $F$-module.

\subsection*{2.2}
\addcontentsline{toc}{subsection}{2.2}

Soit $M$ un motif sur $\Q$, a coefficient dans un corps de nombres $E$. Pour chaque
nombre premier $l$, la realisation $l$-adique $H_l(M)$ de $M$ est un module sur le complete
$l$-adique $E_l$ de $E$. Ce complete est le produit des completes $E_\lambda$ pour $\lambda$
ideal premier au-dessus de $l$, d'ou une decomposition de $H_l(M)$ en un produit de
$E_\lambda$-modules $H_\lambda(M)$.

On conjecture pour les $H_\lambda(M)$ une compatibilite analogue a 1.2.1; si elle est
verifiee, on peut pour chaque plongement complexe $\sigma$ de $E$ definir une serie de
Dirichlet a coefficients dans $\sigma E$, convergente pour $\Re s$ assez grand:
\[
L(\sigma, M, s) = \prod_p L_p(\sigma, M, s), \quad \text{ou}
\]
\begin{equation}\label{eq:2.2.1}
L_p(\sigma, M, s) =_{dfn} Z_p(M, p^{-s}) \text{ avec } Z_p \in E(t) \subset E_\lambda(t)
\text{ donne par}
\end{equation}
$Z_{p\lambda}(M, t) = \det(1 - F_p t, H_\lambda(M)^{I_p})^{-1}$ pour $\lambda \nmid p$.

Pour $\sigma$ variable, ces series de Dirichlet se deduisent les unes des autres par con-
jugaison des coefficients. Nous regarderons le systeme des $L(\sigma, M, s)$ comme une
fonction $L^*(M, s)$ a valeurs dans la $\C$-algebre $E \otimes_\Q \C$, identifiee a $\C^{\Hom(E,\C)}$ par
\begin{equation}\label{eq:2.2.2}
E \otimes_\Q \C \xrightarrow{\sim} \C^{\Hom(E,\C)}: e \otimes z \mapsto (z \cdot \sigma(e))_\sigma.
\end{equation}
Cette fonction peut aussi etre definie directement par un produit eulerien. On espere
comme en 1.1 qu'elle admet un prolongement analytique en $s$.

Il y a lieu de completer le produit eulerien $L(\sigma, M, s)$ par un facteur a l'infini
$L_\infty(\sigma, M, s)$ dependant de la realisation de Hodge de $M$. Posant
$\Lambda(\sigma, M, s) = L_\infty(\sigma, M, s) L(\sigma, M, s)$, l'equation fonctionnelle conjecturale
des fonctions $L$ s'ecrit
\begin{equation}\label{eq:2.2.3}
\Lambda(\sigma, M, s) = \varepsilon(\sigma, M, s) \Lambda(\sigma, \check{M}, 1-s),
\end{equation}
o\`u $\varepsilon(\sigma, M, s)$, comme fonction de $s$, est le produit d'une constante par une exponentielle.
Les definitions de $L_\infty$ et $\varepsilon$ sont rappelees en 5.2. Comme ci-dessus, on regardera le
systeme des $\Lambda$ et celui des $\varepsilon$, pour $\sigma$ variable, comme des fonctions
$\Lambda^*$ et $\varepsilon^*$ a valeurs dans $E \otimes_\Q \C$.

Il resulte de 2.5 ci-dessous que $L_\infty(\sigma, M, s)$ est independant de $\sigma$, et que la fonction
$L_\infty$ pour le motif $R_{E/\Q} M$ deduit de $M$ par restriction du corps des coefficients
(2.1) est la puissance $[E:\Q]$-ieme de $L_\infty(\sigma, M, s)$. Ceci justifie la

\begin{proposition-definition}\label{prop:2.3}
Soit $M$ un motif sur $\Q$ a coefficient dans $E$. Un entier $n$ est critique pour $M$ si les
conditions equivalentes suivantes sont verifiees

\textup{(i)} l'entier $n$ est critique pour $R_{E/\Q} M$;

\textup{(ii)} $L_\infty(\sigma, M, s)$ et $L_\infty(\sigma, \check{M}, s)$ n'ont pas de pole en $s = n$.

On dit que $M$ est critique si $0$ est critique pour $M$.
\end{proposition-definition}

Notre but est de conjecturer la valeur de $L^*(M) =_{dfn} L^*(M, 0)$, pour $M$ critique,
a multiplication par un element de $E$ pres. En d'autres termes, il s'agit de conjecturer
simultanement les valeurs des $L(\sigma, M) =_{dfn} L(\sigma, M, 0)$, a multiplication pres par
un systeme de nombres $\sigma(e)$, $e \in E$.

\subsection*{2.4}
\addcontentsline{toc}{subsection}{2.4}

La realisation $H_B(M)$ de $M$ en cohomologie rationnelle est munie d'une structure
de $E$-espace vectoriel. Sa dimension est le rang sur $E$ de $M$. L'involution $F_\infty$
est $E$-lineaire; les parties $+$ et $-$ sont donc des $E$-sous-espaces vectoriels. On note
$d^+(M)$ et $d^-(M)$ leur dimension.

Le complexifie de $H_B(M)$ est un $E \otimes_\Q \C$-module libre. Identifiant
$E \otimes_\Q \C$ a $\C^{\Hom(E,\C)}$ (2.2.2), on en deduit une decomposition
\[
H_B(M) \otimes_\Q \C = \bigoplus_\sigma H_B(\sigma, M),
\]
avec
\[
H_B(\sigma, M) = (H_B(M) \otimes_\Q \C) \otimes_{E \otimes \C, \sigma} \C,
\]
soit
\[
H_B(\sigma, M) = H_B(M) \otimes_{E, \sigma} \C.
\]
Des $H_B(\sigma, M)$, la decomposition de Hodge etant stable par $E$, chaque
$H_B(\sigma, M)$ herite d'une decomposition de Hodge
$H_B(\sigma, M) = \bigoplus H^{p,q}_B(\sigma, M)$. L'involution $F_\infty$ permute $H^{p,q}$ et $H^{q,p}$,
en particulier stabilise $H^{p,p}_B$ qu'elle decoupe en des parties $+$ et $-$. On note
$h^{p,q}(\sigma, M)$ la dimension de $H^{p,q}_B(\sigma, M)$, et $h^{p,p\pm}(\sigma, M)$ celle de
$H^{p,p\pm}_B(\sigma, M)$. La proposition suivante permet d'omettre $\sigma$ de la notation.

\begin{proposition}\label{prop:2.5}
Les nombres $h^{p,q}(\sigma, M)$ et $h^{p,p\pm}(\sigma, M)$ sont independants de $\sigma$.
\end{proposition}

On peut supposer, et on suppose, que $M$ est homogene. Dans ce cas, $H_B(M)$
s'identifie au complexifie du $E$-espace vectoriel $\bigoplus_{p+q=w} H^{p,q}_B(M)$: c'est un
$E \otimes_\Q \C$-module libre, et la premiere assertion en resulte. Pour la seconde, on
observe que $h^{p,p+}(\sigma, M)$ est l'exces de $d^+(M)$ sur $\sum_{p \neq q} h^{p,q}(\sigma, M)$.

\subsection*{2.6}
\addcontentsline{toc}{subsection}{2.6}

Soit $M$ un motif sur $\Q$ a coefficient dans $E$. Dans la fin de ce paragraphe, sauf
mention expresse du contraire, nous supposons que $R_{E/\Q} M$ verifie les hypotheses
de 1.7. Les espaces $F^\pm$ et $H^\pm_B$ sont cette fois des espaces vectoriels sur $E$. Les
isomorphismes (0.4.1) et (1.7.1)
\begin{align}
I^\pm: H^\pm_B(M) \otimes_\Q \C &\xrightarrow{\sim} H^\pm_{DR}(M) \otimes_\Q \C, \label{eq:2.6.1}\\
I: H_B(M) \otimes_\Q \C &\xrightarrow{\sim} H_{DR}(M) \otimes_\Q \C, \label{eq:2.6.2}
\end{align}
sont des isomorphismes de $E \otimes_\Q \C$-modules entre complexifies d'espaces vectoriels
sur $E$. On pose
\[
c^\pm(M) =_{dfn} \det(I^\pm) \in (E \otimes_\Q \C)^*, \quad \delta(M) =_{dfn} \det(I) \in (E \otimes_\Q \C)^*,
\]
le determinant etant calcule dans des bases $E$-rationnelles de $H^\pm_B(M)$ et
$H^\pm_{DR}(M)$ (resp.~$H_B(M)$ et $H_{DR}(M)$). La definition de $\delta(M)$ ne requiert pas les
hypotheses faites sur $M$. A multiplication par un element de $E^*$ pres, ces nombres ne
dependent que de $M$. Il resulte a nouveau de 1.6 que $c^+(M)$, $i^{d^-(M)} c^-(M)$ et
$i^{d^-(M)} \delta(M)$ sont dans $(E \otimes_\Q \R)^*$.

\begin{conjecture}\label{conj:2.7}
Posons $\mathfrak{F}(M) = -w/2$. Si $M$ est critique,

\textup{(i)} $L(\sigma, M, s)$ n'a jamais de pole en $s = 0$, et ne peut s'annuler en $s = 0$ que
pour $\mathfrak{F}(M) = \tfrac{1}{2}$;

\textup{(ii)} La multiplicite du zero de $L(\sigma, M, s)$ en $s = 0$ est independante de $\sigma$.
\end{conjecture}

Pour (i), je renvoie a la discussion a la fin de 1.3. Que (ii) soit raisonnable m'a ete
suggere par B. Gross.

\begin{conjecture}\label{conj:2.8}
Pour $M$ critique et $L(\sigma, M) \neq 0$, $L^*(M)$ est le produit de $c^+(M)$ par un
element de $E^*$.
\end{conjecture}

\begin{remark}
Un motif $M$ sur un corps de nombres $k$, a coefficient dans $E$, definit aussi une
fonction $L^*(M, s)$. Ces fonctions sont couvertes par notre conjecture, vu l'identite
$L^*(M, s) = L^*(R_{k/\Q} M, s)$, ou $R_{k/\Q}$ est la restriction des scalaires de $k$ a $\Q$
(appliquer (0.5.3)).

Les fonctions $L(\sigma, M, s)$ sont des produits euleriens, indexes par les places finies
de $k$. Il y a lieu de les completer par les facteurs a l'infini $L_v(\sigma, M, s)$, indexes par
les places a l'infini, dont la definition est rappelee en 5.2. Ils dependent en general de
$\sigma$. Sans $L_\infty(\sigma, R_{k/\Q} M, s)$, produit sur $v \mid \infty$ des $L_v(\sigma, M, s)$, est
independant de $\sigma$.
\end{remark}

\begin{remark}
Soient $F$ une extension de $E$, $\sigma$ le morphisme structural de $E \to F$ et $\iota_\C$ son
complexifie: $E \otimes_\Q \C \hookrightarrow F \otimes_\Q \C$. On a $L^*(M \otimes_E F, s) = \iota_\C L^*(M, s)$
et $c^\pm(M \otimes_E F) = \iota_\C c^\pm(M)$. La conjecture est donc compatible a l'extension du
corps des coefficients. Pour $F$ une extension galoisienne de $E$, de groupe de Galois
$G$, le Theoreme 90 de Hilbert $H^1(G, F^*) = 0$ assure que
$((F \otimes_\Q \C)^*/F^*)^G = (E \otimes_\Q \C)^*/E^*$: la conjecture est invariante par
extension du corps des coefficients.
\end{remark}

\begin{remark}
Si $E$ est une extension de $F$, on a $L^*(R_{E/F} M, s) = N_{E/F} L^*(M, s)$. Les periodes
$c^\pm$ verifiant la meme identite, la conjecture est compatible a la restriction des coefficients.
\end{remark}

\begin{remark}
Soit $D$ une algebre a division de rang $d^2$ sur $E$. Un motif $M$ sur $\Q$, muni d'une
structure de $D$-module, et de rang $n$ sur $D$, definit une serie de Dirichlet a coefficients
dans $E$ dont les facteurs euleriens sont presque tous de degre $nd$: pour $\lambda$ une place
finie de $E$, $H_\lambda(M)$ est un module libre sur le complete $D_\lambda = D \otimes_E E_\lambda$,
on reprend la definition (2.2.1) en posant
\[
Z_p(M, t) = \detred(1 - F_p t, H_\lambda(M)^{I_p})^{-1}
\]
(si $D_\lambda$ est une algebre de matrices sur $E_\lambda$, et $e$ un idempotent indecomposable, le
determinant reduit d'un endomorphisme $A$ d'un $D_\lambda$-module $H$ est le determinant,
calcule sur $E_\lambda$, de la restriction de $A$ a $eH$; la definition dans le cas general procede
par descente).

La liberte que nous donne 2.10 d'etendre le corps des coefficients met ces fonctions $L$
elles aussi sous le chapeau 2.8: choisissant une extension $\xi: E \hookrightarrow F$ de $E$
qui neutralise $D$, et un idempotent indecomposable $e$ de $D \otimes_E F$, on a
$e L^*(M, s) = L^*(e(M \otimes_E F), s)$.

On peut aussi definir $c^\pm(M)$ directement dans ce cadre. C'est ce qui est explique
en 2.13 ci-dessous.
\end{remark}

La fin de ce paragraphe est inutile pour la suite.

\subsection*{2.13}
\addcontentsline{toc}{subsection}{2.13}

Soit $D$ une algebre simple sur un corps $E$ (voire une algebre d'Azumaya sur un
anneau...) Pour le formalisme tensoriel, il est commode de regarder les $D$-modules
comme de ``faux $E$-espaces vectoriels'':

(a) Se donner un espace vectoriel $V$ sur $E$ revient a se donner, pour toute extension
etale $F$ de $E$, un $F$-espace vectoriel $V_F$, et des isomorphismes compatibles
$V_G = V_F \otimes_F G$ pour $G$ une extension de $F$. On prend $V_F = V \otimes_E F$. Par descente,
il suffit de ne se donner les $V_F$ que pour $F$ assez grand.

(b) Soit $W$ un $D$-module. Pour toute extension etale $F$ de $E$, et tout
$F$-isomorphisme $D \otimes_E F \xrightarrow{\sim} \End_F(L)$, avec $L$ libre, posons
$W_{L,F} = \Hom_{D \otimes F}(L, W \otimes_E F)$ (les produits tensoriels sont sur $E$). On a un
isomorphisme de $D \otimes F$-modules: $W \otimes_E F = L \otimes_F W_{L,F}$.

Si $F$ est assez grand pour neutraliser $D$, $L$ est unique: a isomorphisme non unique
pres; la non-unicite est due aux homotheties, qui agissent trivialement sur $\End(L)$.
C'est pourquoi la donnee des $W_{L,F}$ n'est pas du type (a); c'est un ``faux espace
vectoriel sur $E$''.

Soient $(W_i)$ une famille de $D$-modules, et $T$ une operation tensorielle. Si les
homotheties de $L$ agissent trivialement sur $T(W_i \otimes L)$, le $F$-espace vectoriel
$T(W_i \otimes L)$ est independant du choix de $L$ et on obtient un systeme du type (a), d'ou un espace
vectoriel $T(W_i)$ sur $E$.

\begin{example}
Si $W'$ et $W''$ sont deux $D$-modules de rang $n \cdot [D:\!E]^{1/2}$ sur $E$, on peut
prendre $T = \Hom(\bigwedge^n W'', \bigwedge^n W' \otimes L)$; on obtient un espace vectoriel
$E(W', W'')$ de rang 1 sur $E$, et tout homomorphisme $f: W' \to W''$ a un determinant
reduit $\detred(f) \in E(W', W'')$.
\end{example}

Pour definir $c^+(M)$, on applique cette construction aux $D$-modules $H^+_B(M)$ et
$H^+_{DR}(M)$. Posons $\mathcal{E} = \mathcal{E}(H^+_B(M), H^+_{DR}(M))$. Le determinant reduit de
l'isomorphisme de $D \otimes_\Q \C$-modules $I^+: H^+_B(M)_\C \xrightarrow{\sim} H^+_{DR}(M)_\C$
est dans le $E \otimes \C$-module libre de rang 1 $\mathcal{E} \otimes \C$. On pose
$\detred(I^+) = c^+(M) \cdot e$ pour $e$ une base de $\mathcal{E}$.

\section*{3. Exemple: la fonction $\zeta$}
\addcontentsline{toc}{section}{3. Exemple: la fonction zeta}

\subsection*{3.1}
\addcontentsline{toc}{subsection}{3.1}

Pour comprendre les diverses realisations du motif de Tate $\Z(1)$, le plus simple est
de l'ecrire $\Z(1) = H^1(\mathbb{G}_m)$. Le groupe multiplicatif n'etant pas une variete
projective, ceci ne rentre pas dans le cadre de Grothendieck, qui demande qu'on
definisse plutot $\Z(1)$ comme le dual du facteur direct $H^2(\PP^1)$ de $H^*(\PP^1)$.

La realisation en $\Z_l$-cohomologie de $\Z(1)$ est le module de Tate $T_l(\mathbb{G}_m)$ de
$\mathbb{G}_m$
\[
T_l(\mathbb{G}_m) = \varprojlim \mu_{l^n}.
\]
En realisation de Hodge, on a $H_B(\Z(1)) = H^1(\C^*)$ isomorphe a $\Z(2\pi i)$ ($\Q$
plutot, en homologie rationnelle). Ce groupe est purement de type $(-1,-1)$ et
$F_\infty = -1$.

En realisation de de Rham, $H_{DR}(\Z(1))$ est le dual de $H^1_{DR}(\mathbb{G}_m)$, isomorphe a
$\Q$, de generateur la classe de $dz/z$. L'unique periode de $H^1(\mathbb{G}_m)$ est
\begin{equation}\label{eq:3.1.1}
\int_{S^1} \frac{dz}{z} = 2\pi i.
\end{equation}

Sur $\Z_l(1)$, le Frobenius arithmetique $\varphi_p$ ($p \neq l$) agit par multiplication par $p$. Le
Frobenius geometrique agit donc par multiplication par $p^{-1}$; ceci justifie l'identite
citee en 1.3
\begin{equation}\label{eq:3.1.2}
L(M(n), s) = L(M, n+s).
\end{equation}

Puisque $F_\infty$ agit sur $(-1)^n$ sur $H_B(\Z(n))$, $H^\alpha_B(\Z(n))$ est nul pour
$\alpha = -(-1)^n$, et $c^+(\Z(n)) = 1$. Pour $\alpha = (-1)^n$, d'apres (3.1.1),
$c^-(\Z(n)) = \delta(\Z(n)) = (2\pi i)^n$:
\begin{equation}\label{eq:3.1.3}
\begin{aligned}
c^+(\Z(n)) &= (2\pi i)^n &&\text{pour } \alpha = (-1)^n,\\
c^+(\Z(n)) &= 1, \quad \delta(\Z(n)) = (2\pi i)^n &&\text{pour } \alpha = -(-1)^n.
\end{aligned}
\end{equation}

\subsection*{3.2}
\addcontentsline{toc}{subsection}{3.2}

La fonction $\zeta(s)$ est la fonction $L$ attachee au motif unite $\Z(0) = H^*(\text{Point})$.
Les entiers critiques pour $\Z(0)$ sont les entiers pairs $> 0$, et les entiers impairs
$\leq 0$. A cause du pole de $\zeta(s)$ en $s = 1$, $0$ n'est pas un zero trivial et il serait
raisonnable de definir les entiers critiques comme incluant $0$. Les equations (3.1.3)
et les valeurs connues de $\zeta(n) = L(\Z(n))$ pour $n$ critique verifient 1.8:
$\zeta(n)$ est rationnel pour $n$ impair $\leq 0$, et un multiple rationnel de $(2\pi i)^{|n|/2}$
pour $n$ pair $> 0$.


\section*{4. Compatibilite a la conjecture de Birch et Swinnerton-Dyer}
\addcontentsline{toc}{section}{4. Compatibilite a la conjecture de Birch et Swinnerton-Dyer}

\subsection*{4.1}
\addcontentsline{toc}{subsection}{4.1}

Soient $A$ une variete abelienne sur $\Q$, et $d$ sa dimension. La conjecture de
Birch et Swinnerton-Dyer [15] affirme notamment:

(a) $L(H^1(A), 1)$ est non nul si et seulement si $A(\Q)$ est fini.

(b) Soit $\omega$ un generateur de $H^0(A, \Omega^d)$. Alors, $L(H^1(A), 1)$ est le produit de
$\int_{A(\R)} |\omega|$ par un nombre rationnel.

Le motif $H^1(A)(1)$ est isomorphe au dual $H^1(A)^\vee$ de $H^1(A)$: ceci traduit l'existence
d'une polarisation, autodualite de $H^1(A)$ a valeurs dans $\Z(-1)$. D'apres 1.7,
$c^+(H^1(A)(1))$ se calcule donc comme suit: si $\omega_1, \ldots, \omega_d$ est une base de
$H^0(A, \Omega^1) = F^+ H^1_{DR}(A)$, et $c_1, \ldots, c_d$ une base de $H_1(A(\C), \Q)^+$, on a
\begin{equation}\label{eq:4.1.1}
c^+(H^1(A)(1)) = \det(\langle \omega_i, c_j \rangle).
\end{equation}

Designant encore par $c_i$ un cycle representatif, on a $\langle \omega_i, c_j \rangle = \int_{c_j} \omega_i$.
Prenons pour base $(c_i)$ une base sur $\Z$ de $H_1(A(\R)^0, \Z) \subset H_1(A(\C), \Z)$. Le
produit de Pontryagin des $c_i$ est represente par le $d$-cycle $A(\R)^0$ dans $A(\C)$, pour une
orientation convenable, et, si $\omega$ est le produit exterieur des $\omega_i$, le determinant
(4.1.1) est l'integrale $\int_{A(\R)^0} \omega$. On a
\[
\int_{A(\R)^0} \omega = [A(\R): A(\R)^0]^{-1} \int_{A(\R)} |\omega|,
\]
et 1.8 pour $H^1(A)(1)$ equivaut donc a 4.1(b) ci-dessus.

\subsection*{4.2}
\addcontentsline{toc}{subsection}{4.2}

La conjecture de Birch et Swinnerton-Dyer donne la valeur exacte de
$L(H^1(A), 1)$; la description du facteur rationnel en 4.1(b) a partir du motif $H^1(A)(1)$
a la forme suivante:

(a) La donnee de $M = H^1(A)(1)$ equivaut a celle de $A$ a isogenie pres. Il faut
commencer par choisir $A$. Cela revient a choisir un reseau entier dans $H_B(M)$, dont
les $l$-adifies soient stables sous l'action de $\Gal(\overline{\Q}/\Q)$.

(b) On choisit alors $\omega$, par exemple comme produit exterieur des elements d'une
base d'un reseau entier dans $H^0_{DR}(M)$. Cela determine la periode, $c^+(M)$, et, pour
chaque nombre premier $p$, un facteur rationnel $c_p(M)$. Les $c_p(M)$ sont presque tous
egaux a 1, et la formule du produit assure que $c^+(M) \cdot \prod_p c_p(M)$ est independant
de $\omega$.

(c) Un autre nombre rationnel, $h(M)$, est defini en terme d'invariants cohomolo-
giques de $A$; le nombre rationnel cherche est $h(M) \cdot \prod_p c_p(M)^{-1}$.

L'invariance de la conjecture par isogenie est par ailleurs un theoreme non trivial.

\subsection*{4.3}
\addcontentsline{toc}{subsection}{4.3}

Pour generaliser 4.1(a) a un motif de poids $-1$ $M$ quelconque, il faudrait
disposer d'un analogue de $A(\Q)$. Le groupe $A(\Q)$ peut s'interpreter comme le
groupe des extensions de $\Z(0)$ par $H^1(A)$, dans la categorie des 1-motifs [6, \S 10] sur
$\Q$. Ceci suggere de considerer le groupe des extensions de $\Z(0)$ par $M$, dans une
categorie de motifs mixtes, paralleles aux structures de Hodge mixtes, mais on ne
dispose meme pas d'une definition conjecturale d'une telle categorie!

J'observerai seulement que, dans toutes les theories cohomologiques usuelles, un
cycle $Y$ de dimension $d$, cohomologue a zero, sur une variete algebrique propre et
lisse $X$ determine un torseur sous $H^{2d-1}(X)(d)$: on dispose d'une suite exacte de
cohomologie
\[
0 \to H^{2d-1}(X)(d) \to H^{2d-1}(X - Y)(d) \to H^{2d}_Y(X)(d) \to H^{2d}(X)(d)
\]
$Y$ definit une classe de cohomologie $\text{cl}(Y) \in H^{2d}_Y(X)(d)$, d'image nulle dans
$H^{2d}(X)(d)$, et on prend $\partial \text{cl}(Y)$. Cette construction correspond a celle qui, a un
diviseur de degre 0 sur une courbe, associe un point de sa jacobienne.


\section*{5. Compatibilite a l'equation fonctionnelle}
\addcontentsline{toc}{section}{5. Compatibilite a l'equation fonctionnelle}

Soit $E$ un corps de nombres. Dans $(E \otimes_\Q \C)^*$, nous noterons $\sim$ la relation
d'equivalence definie par le sous-groupe $E^*$.

\begin{proposition}\label{prop:5.1}
Soit $M$ un motif sur $\Q$, a coefficient dans $E$. On suppose verifiees les hypotheses
de 1.7. On a alors
\[
c^+(M) \sim (2\pi i)^{d^-(M)} \cdot \mathfrak{F}(M) \cdot c^+(\check{M}(1)).
\]
\end{proposition}

Pour $L$ un module libre de rang $n$ sur un anneau commutatif $A$ ($A$ sera $E$, ou
$E \otimes_\Q \C$), posons $\det L = \bigwedge^n L$. On etend par localisation cette definition au cas ou
$L$ est seulement projectif de type fini (cette generalisation n'est pas indispensable a la
preuve de 5.1). On a des isomorphismes canoniques
\begin{align}
\det(L^\vee) &= (\det L)^{-1}, \label{eq:5.1.1}\\
\intertext{et, pour tout facteur direct $P$ et $L$,}
\det(L) &= \det(P) \cdot \det(L/P) \label{eq:5.1.2}
\end{align}
(on a designe par $\cdot$ un produit tensoriel, et par $^{-1}$ un dual). Il y a ici des problemes
de signe, qu'on resoud au mieux en considerant $\det(L)$ comme un module gradue
inversible, place en degre le rang de $L$. Nos resultats finaux etant modulo $\sim$, nous
ne nous en inquieterons pas.

Pour $X$ et $Y$ de meme rang, posons $\mathcal{B}(X, Y) = \Hom(\det X, \det Y) =
(\det X)^{-1} \cdot \det(Y)$. Le determinant de $f: X \to Y$ est dans $\mathcal{B}(X, Y)$. On deduit de (5.1.1)
et (5.1.2) des isomorphismes
\begin{align}
\mathcal{B}(X, Y) &= \mathcal{B}(Y^\vee, X^\vee), \label{eq:5.1.3}\\
\intertext{et, pour $F \subset X$ et $G \subset Y$,}
\mathcal{B}(X, Y) &= \mathcal{B}(F, Y/G) \cdot \mathcal{B}(G, X/F)^{(-1)^{\dim G}}. \label{eq:5.1.4}
\end{align}

\begin{lemma}\label{lem:5.1.5}
Via l'isomorphisme (5.1.3), on a $\det(f) = \det(\check{f})$.
\end{lemma}

\begin{lemma}\label{lem:5.1.6}
Soient $F$ et $G$ des facteurs directs de $X$ et $Y$. On suppose que l'isomorphisme
$f: X \to Y$ induit des isomorphismes $f_F: F \to Y/G$ et $f_G: G \to X/F$. Via
l'isomorphisme (5.1.4), on a $\det(f) = \det(f_F) \cdot \det(f_G)^{(-1)^{\dim G}}$.
\end{lemma}

La verification de ces lemmes est laissee au lecteur. Pour 5.1.6, notons seulement
que, pour $G' = f^{-1}(G)$ et $F_1 = f(F)$, on a $X = F \oplus G'$, $Y = G \oplus F_1$, et que $f$
echange $F$ et $F_1$, $G'$ et $G$.

Appliquons le Lemme 5.1.6 aux complexifies de $H^\pm_B(M) \subset H_B(M)$ et de
$F^\pm \subset H_{DR}(M)$. On trouve que, via l'isomorphisme (5.1.4), le determinant de
$I: H_B(M) \otimes \C \to H_{DR}(M) \otimes \C$ est le produit du determinant de
$I^+: H^+_B(M) \otimes \C \to (H_{DR}(M)/F^-) \otimes \C$ et de l'inverse du determinant du morphisme
induit par l'inverse de $I$:
\[
I^-: F^- \otimes \C \to (H_B(M)/H^+_B(M)) \otimes \C.
\]
Le morphisme $I^-$ est le transpose du morphisme $I^+$ pour le motif $\check{M}$ dual de $M$.
Appliquant 5.1.5, et prenant des bases sur $E$ des espaces $\mathcal{B}(X, Y)$ en jeu, on obtient
finalement que
\begin{equation}\label{eq:5.1.7}
\delta(M) \sim c^+(M) \cdot c^-({\check{M}})^{-1}.
\end{equation}
On en deduit 5.1 en appliquant au motif $\check{M}$ la formule suivante, consequence
de 3.1:
\begin{equation}\label{eq:5.1.8}
c^\pm(M) = (2\pi i)^{-d^\mp(M) \cdot \mathfrak{F}(M)} c^\pm(M(n)).
\end{equation}
Notons aussi, pour usage ulterieur, la formule analogue
\begin{equation}\label{eq:5.1.9}
\delta(M) = (2\pi i)^{-d^-(M) \cdot \mathfrak{F}(M)} \delta(M(n)).
\end{equation}

\subsection*{5.2}
\addcontentsline{toc}{subsection}{5.2}

Rappelons la forme exacte de l'equation fonctionnelle conjecturale des fonctions
$L$ des motifs ([12], [4]). Nous nous placerons dans le cas general d'un motif sur un
corps de nombres $k$, a coefficient dans un corps $E$ muni d'un plongement complexe
$\sigma$ (cf.~2.9). La forme generale d'abord:

(a) Pour chaque place $v$ de $k$, on definit un facteur local $L_v(\sigma, M, s)$. Pour $v$ fini,
la definition de $L_v(\sigma, M, s) = Z_v(M, N_v^{-s})$ ($Z_v(M, t) \in E(t)$) depend d'une hypothese
de compatibilite analogue a (1.2.1). On a $Z_v(M, t) = \det(1 - F_v t, H_l(M)^{I_v})^{-1}$.
Pour $v$ infini, induit par un plongement complexe $\tau$, on obtient $L_v(\sigma, M, s)$ en
decomposant $H_\tau(M) \otimes_{E,\sigma} \C$ en somme directe de sous-espaces minimaux stables par
les projecteurs qui donnent la decomposition de Hodge, et par $F_\infty$ pour $v$ reel, en
associant a chacun le ``facteur $L$'' de la Table (5.3), et en prenant leur produit. Pour
$v$ complexe, les sous-espaces minimaux sont de dimension un, d'un type $(p, q)$.
Pour $v$ reel, ils sont soit de dimension 2, de type $\{(p,q), (q,p)\}$, $p \neq q$, soit de
dimension 1, de type $(p,p)$, avec $F_\infty = \pm 1$. On note $\Lambda(\sigma, M, s)$ le produit des
$L_v(\sigma, M, s)$.

(b) Soient $\psi$ un caractere non trivial du groupe $A \otimes_k k / k$ des classes d'adeles de
$k$, et $\psi_v$ ses composantes. Soient aussi, pour chaque place $v$ de $k$, une mesure de
Haar $dx_v$ sur $k_v$. On suppose que, pour presque tout $v$, $dx_v$ donne aux entiers de
$k_v$ la masse 1, et que le produit $\bigotimes dx_v$ des $dx_v$ est la mesure de Tamagawa, donnant
la masse 1 au groupe des classes d'adeles.

On definit des constantes locales $\varepsilon_v(\sigma, M, s, \psi_v, dx_v)$, presque toutes egales a 1 et
toutes, comme fonction de $s$, produit d'une constante par une exponentielle. Soit
$\varepsilon(\sigma, M, s)$ leur produit (independant de $\psi$ et des $dx_v$).

(c) L'equation fonctionnelle conjecturale est
\[
\Lambda(\sigma, M, s) = \varepsilon(\sigma, M, s) \cdot \Lambda(\sigma, \check{M}, 1-s).
\]

Pour definir les $\varepsilon_v$, une hypothese additionnelle de compatibilite entre les $H_l(M)$
est requise. Elle permet d'associer a $v$, $\sigma$, $M$ une classe d'isomorphie de representations
complexes du groupe de Weil $W(\overline{k}_v/k_v)$ pour $v$ infini, du groupe de Weil epaissi
$\mathcal{W}'(\overline{k}_v/k_v)$ pour $v$ fini, et on prend le $\varepsilon$ de [4, 8.12.4], avec $|t|_v = q_v^{-\deg}$.

Pour $v$ infini, ceci revient a decomposer $H_\sigma(M) \otimes_{E,\sigma} \C$ comme en (a), a associer
a chaque sous-espace de la decomposition un facteur $\varepsilon_v$, et a prendre leur produit.
La table des $\varepsilon_v$, pour un choix particulier de $\psi_v$ et de $dx_v$, est donnee en 5.3.

Pour $v$ fini, on commence par restreindre les representations $H_l(M)$ a un groupe de
decomposition $\Gal(\overline{k}_v/k_v) \subset \Gal(\overline{k}/k)$, puis au groupe de Weil $W(\overline{k}_v/k_v)$. Appliquant
[4, 8.3, 8.4], on deduit de $H_B(M)$ une classe d'isomorphie de representations de
$\mathcal{W}'(\overline{k}_v/k_v)$ sur $E_\lambda$. Il est loisible ici, et utile, de la remplacer par sa $F$-semi-simplifiee
[4, 8.6]. On demande que ces representations, pour $\lambda$ variable, soient compatibles,
i.e., que si on etend les scalaires de $E_\lambda$ a $\C$, par $\tilde{\sigma}: E_\lambda \to \C$ prolongeant $\sigma$,
la classe d'isomorphie de la representation obtenue soit independante de $\lambda$ et de $\sigma$.
C'est la classe d'isomorphie cherchee.

Remontons pour le lecteur les renvois internes de [4]. Une representation de $\mathcal{W}'$
est donnee par une representation $\rho$ du groupe de Weil dans $\GL(V)$, et par un
endomorphisme nilpotent $N$ de $V$. En terme de la constante locale [4, 4.1] de $\rho$,
celle de $(\rho, N)$ est donnee par
\[
\varepsilon((\rho, N), s, \psi, dx) = \varepsilon(\rho \otimes |\cdot|^s, \psi, dx) \cdot \det(-F \cdot N_v^{-s}, V^{I_v}/\Ker(N)^{I_v}).
\]

\begin{remark}
La fonction de $s$
\[
\varepsilon((\rho, N), s, \psi, dx) L((\rho, N)^\vee, 1-s) / L((\rho, N), s)
\]
est la meme pour $(\rho, N)$ et pour $(\rho, 0)$. Ceci permet d'enoncer l'equation fonction-
nelle conjecturale des fonctions $L$ en supposant seulement une compatibilite entre
les semi-simplifiees des restrictions des $H_l(M)$ au groupe de decomposition.
\end{remark}

\subsection*{5.3}
\addcontentsline{toc}{subsection}{5.3}

Dans la table ci-dessous, on donne les facteurs locaux, et les constantes, associes
aux divers types de sous-espaces minimaux de la realisation de Hodge. Pour les
constantes, on a suppose $\psi$ et la mesure $dx_v$ choisis comme suit:

$v$ reel: $\psi_v(x) = \exp(2\pi i x)$, mesure $dx_v$;

$v$ complexe: $\psi_v(z) = \exp(2\pi i \Tr_{\C/\R}(z))$, mesure $|dz \wedge d\bar{z}|$,
soit pour $z = x + iy$, $\exp(4\pi i x)$ et $dx \, dy$.

On utilise les notations $\Gamma_\R(s) = \pi^{-s/2} \Gamma(s/2)$, $\Gamma_\C(s) = 2 \cdot (2\pi)^{-s} \Gamma(s)$.

\begin{center}
\begin{tabular}{llll}
\toprule
place & type & facteur $L$ & constante $\varepsilon$ \\
& & & (pour $\psi_v$, $dx_v$ ci-dessus) \\
\midrule
complexe & $(p,q)$ ou $(q,p)$, $p \leq q$ & $\Gamma_\C(s-p)$ & $i^{q-p}$ \\
reelle & $\{(p,q), (q,p)\}$, $p < q$ & $\Gamma_\R(s-p)\Gamma_\R(s-q)$ & $i^{q-p+1}$ \\
& $(p,p)$, $F_\infty = (-1)^{p+\varepsilon}$, $\varepsilon = 0$ ou $1$ & $\Gamma_\R(s+\varepsilon-p)$ & $i^\varepsilon$ \\
\bottomrule
\end{tabular}
\end{center}

Le cas qui nous interesse est celui ou $k = \Q$. On peut dans ce cas prendre
$\psi_\infty(x) = \exp(2\pi i x)$, $\psi_p(x) = \exp(-2\pi i x)$ (via l'isomorphisme $\Q_p/\Z_p =$
partie $p$-primaire de $\Q/\Z$), $dx_\infty =$ mesure de Lebesgue $dx$, $dx_p =$ la mesure de
Haar sur $\Q_p$, donnant a $\Z_p$ la masse 1.

\begin{proposition}\label{prop:5.4}
Si $M$ est critique de poids $w$, on a, modulo un nombre rationnel independant de $\sigma$,
\[
L_\infty(\sigma, M(1)) L_\infty(\sigma, M)^{-1} \sim (2\pi i)^{d^-(M)} \cdot (2\pi)^{-w \cdot d(M)/2}.
\]
\end{proposition}

D'apres 2.5, les $L_\infty(\sigma, \cdot)$ sont independants de $\sigma$. Ceci nous permet de ne
verifier 5.4 que pour $\sigma$ fixe. La formule est compatible a la substitution
$M \mapsto M(1)$: $M(1)$ est de poids $-2-w$, son $\mathfrak{F}$ est $(1+\mathfrak{F}(M))$ et, abregeant
$d(M)$ et $d^\pm(M)$ en $d$, $d^\pm$, on a $d = d^+ + d^-$ et
\[
(-d^- - w \cdot d/2) + (-d^+ - (-2-w) \cdot d/2) = 0.
\]
Ceci permet de ne verifier 5.4 que pour $w \geq -1$.

Pour $s$ entier, on a, modulo $\Q^*$:
\begin{equation}\label{eq:5.4.1}
\begin{aligned}
\Gamma_\R(s) &\sim (2\pi)^{-s/2} &&\text{pour $s$ pair } > 0,\\
\Gamma_\R(s) &\sim (2\pi)^{(1-s)/2} &&\text{pour $s$ impair},\\
\Gamma_\C(s) &\sim (2\pi)^{-s} &&\text{pour $s > 0$.}
\end{aligned}
\end{equation}

Pour $w \geq -1$, la puissance de $2\pi i$ dans la contribution de chaque sous-espace de
$H_\sigma(M) \otimes_\Q \C$, comme en 5.2(a), est donc donnee par

\begin{center}
\begin{tabular}{llll}
\toprule
& $L_\infty(M)$ & $L_\infty(M(1))$ & $L_\infty(M(1)) L_\infty(M)^{-1}$ \\
\midrule
$\{(p,q), (q,p)\}$, $p \leq q$ & $-p$ & $-1-p$ & $-1-w$ \\
$(p,p)$, $p$ pair; $0$, $F_\infty = 1$ & $\tfrac{1}{2}-p$ & $-\tfrac{1}{2}-p$ & $-1$ \\
$(p,p)$, $p$ impair; $0$, $F_\infty = -1$ & $\tfrac{1}{2}-p$ & $\tfrac{3}{2}-p$ & $1$ \\
\bottomrule
\end{tabular}
\end{center}
La proposition en resulte aussitot.

Posons $\det M =_{dfn} \bigwedge^{\rank_E M} M$ (puissance exterieure sur $E$).

\begin{proposition}\label{prop:5.5}
$\varepsilon^*(M) \sim \varepsilon^*(\det M)$.
\end{proposition}

Pour des $dx_v$ choisis comme suggere en 5.4, nous prouverons plus precisement
des equivalences
\begin{equation}\label{eq:5.5.1}
\varepsilon_v(M, \psi_v, dx_v) \sim \varepsilon_v(\det M, \psi_v, dx_v).
\end{equation}
Posant $\eta_v(\sigma, M, \psi_v) = \varepsilon_v(\sigma, M, \psi_v, dx_v) \cdot \varepsilon_v(\sigma, \det M, \psi_v, dx_v)^{-1}$,
ceci equivaut a
\begin{equation}\label{eq:5.5.2}
\eta_v(\sigma, M, \psi_v) = \eta_v(\sigma, M, \psi_v^a)
\end{equation}
pour tout automorphisme $I$ de $\C$.

Il resulte de [4, 5.4] que, si $a \in k_v^*$ est de valeur absolue 1, et qu'on pose
$(\psi_v \cdot a)(x) = \psi_v(ax)$, on a
\begin{equation}\label{eq:5.5.3}
\eta_v(\sigma, M, \psi_v) = \eta_v(\sigma, M, \psi_v \cdot a) \quad \text{(pour } |a|_v = 1\text{).}
\end{equation}
Pour $I$ un automorphisme de $\C$, et $v$ fini, $I \psi_v$ est de cette forme $\psi_v \cdot a$
avec $|a|_v = 1$; de meme, $I \psi_\infty = \psi_\infty \cdot (-1)$.

Pour $v$ fini, la definition de $\varepsilon_v$ est purement algebrique, d'ou
$\eta_v(\sigma, M, \psi_v) = \eta_v(\sigma, M, I \psi_v)$ et (5.5.2) resulte de (5.5.3).

Pour $v = \infty$, on a encore $\varepsilon_\infty(\sigma, M, I \psi_\infty) = \eta_\infty(\sigma, M, \psi_\infty) =
\eta_\infty(\sigma, M, \psi_\infty)$; si on prend $\psi_\infty$ comme suggere en 5.3, $I \psi_\infty$ est une
puissance de $i$ independante de $\sigma$, d'ou $\eta_\infty(\sigma, M, I \psi_\infty) = i^k$, independant de
$\sigma$, ce qui verifie (5.5.2).

\begin{theorem}\label{thm:5.6}
Modulo la Conjecture 6.6 sur la nature des motifs de rang 1, la Conjecture 2.8
est compatible a l'equation fonctionnelle conjecturale des fonctions $L$: on a
\[
L^*(M) \cdot c^+(M) \sim \varepsilon^*(M) \cdot L^*_\infty(\check{M}(1)) \cdot c^+(\check{M}(1)).
\]
\end{theorem}

D'apres 5.1, 5.3 et 5.5, cette formule equivaut a:
\[
(2\pi i)^{d^-(M) \cdot \mathfrak{F}(M)} \cdot \delta(M) \sim (2\pi i)^{d^-(M) \cdot \mathfrak{F}(M)} \cdot (2\pi)^{-w \cdot d(M)/2} \cdot \varepsilon^*(\det M).
\]
Posons $D = \det M$ et $a = d^-(D)$. On a $\delta(M) = \delta(D)$, $d^-(M) \equiv a \pmod 2$,
et $w \cdot d(M)$ est le poids $w(D)$ de $D$, de sorte que la formule equivaut encore a
\begin{equation}\label{eq:5.6.1}
\varepsilon^*(D) \sim (2\pi i)^{a + w(D)/2} \cdot i^a \cdot \delta(D).
\end{equation}
Nous prouverons en 6.5 que (5.6.1) vaut pour une classe de motifs de rang 1 qui,
conjecturalement (6.6), les englobe tous.

\section*{6. Exemple: les fonctions $L$ d'Artin}
\addcontentsline{toc}{section}{6. Exemple: les fonctions L d'Artin}

\begin{definition}\label{def:6.1}
La categorie des motifs d'Artin est l'enveloppe karoubienne de la duale de la
categorie d'objets les varietes sur $\Q$ de dimension 0, de morphismes les
correspondances definies sur $\Q$.
\end{definition}

Par definition, chaque variete de dimension 0, $X$, definit un motif d'Artin
$H(X)$, le foncteur $H$ est un foncteur contravariant pleinement fidele,
\[
H: (\text{varietes de dim } 0, \text{correspondances}) \to (\text{motifs d'Artin})
\]
et tout motif d'Artin est facteur direct d'un $H(X)$.

\subsection*{6.2}
\addcontentsline{toc}{subsection}{6.2}

Explicitons cette definition. Soient $\overline{\Q}$ une cloture algebrique de $\Q$, et $G$ le
groupe de Galois $\Gal(\overline{\Q}/\Q)$. Une variete de dimension 0 est le spectre d'un produit
fini $A$ de corps de nombres, et la theorie de Galois (sous la forme que l'a donnee
Grothendieck) dit que le foncteur
\[
X = \Spec(A) \mapsto X(\overline{\Q}) = \Hom(A, \overline{\Q}):
\]
$($categorie des varietes de dimension 0, sur $\Q$, et des morphismes de schemas$)$
$\to ($categorie des ensembles finis munis d'une action continue de $G)$
est une equivalence de categorie. Le foncteur inverse est $I \mapsto$ spectre de l'anneau
des fonctions $G$-invariantes de $I$ dans $\overline{\Q}$.

Une correspondance d'une variete de dimension 0, $X$ dans une autre, $Y$, est une
combinaison lineaire formelle a coefficients dans $\Q$ de composantes connexes de
$X \times Y$. L'application
\[
\sum a_i Z_i \mapsto \sum a_i (\text{fonction caracteristique de } Z_i(\overline{\Q}) \subset (X \times Y)(\overline{\Q}))
\]
identifie correspondances et fonctions $G$-invariantes, a valeurs rationnelles, sur
$(X \times Y)(\overline{\Q}) = X(\overline{\Q}) \times Y(\overline{\Q})$, et la composition des correspondances au produit
matriciel.

Notons $H$ le foncteur contravariant $X \mapsto$ espace vectoriel $\Q^{X(\overline{\Q})}$, muni de l'action
naturelle de $G$; $($correspondance $F: X \to Y) \mapsto$ le morphisme
$F^*: \Q^{X(\overline{\Q})} \to \Q^{Y(\overline{\Q})}$ de matrice $^tF$. Il est pleinement fidele, et identifie la categorie
des motifs d'Artin a celle des representations rationnelles de $G$.

\subsection*{6.3}
\addcontentsline{toc}{subsection}{6.3}

Dans ce modele, si $\overline{\Q}$ est la cloture algebrique de $\Q$ dans $\C$, le foncteur
``realisation de Betti'' $H_B$ est le foncteur ``espace vectoriel sous-jacent''. La structure
de Hodge est purement de type $(0,0)$, et l'involution $F_\infty$ est l'action de la
conjugaison complexe $F_\infty \in G$. On a en effet un isomorphisme, fonctoriel pour les
correspondances,
\[
H^*(X(\C), \Q) = \Q^{X(\C)} = \Q^{X(\overline{\Q})} = H(X).
\]

Le foncteur ``realisation $l$-adique'' $H_l$ est le foncteur $H_l(V) = V \otimes_\Q \Q_l$.
On a en effet un isomorphisme, fonctoriel pour les correspondances,
\[
H^*(X(\overline{\Q}), \Q_l) = \Q_l^{X(\overline{\Q})} = \Q^{X(\overline{\Q})} \otimes_\Q \Q_l.
\]

Calculons de meme la realisation de de Rham. Pour $X = \Spec(A)$, on a
$H_{DR}(X) = A$. Ecrivant que $A = \Hom_G(X(\overline{\Q}), \overline{\Q}) = (\Q^{X(\overline{\Q})} \otimes \overline{\Q})^G$,
on obtient que
\begin{equation}\label{eq:6.3.1}
H_{DR}(V) = (V \otimes_\Q \overline{\Q})^G.
\end{equation}

La formule (6.3.1) realise $H_{DR}(V)$ comme un sous-espace de $V \otimes_\Q \overline{\Q}$. Ce sous-
espace est une $\Q$-structure: on a $(V \otimes_\Q \overline{\Q})^G \otimes_\Q \overline{\Q} \xrightarrow{\sim} V \otimes_\Q \overline{\Q}$.
Apres extension des scalaires a $\overline{\Q}$, $H_B(V)$ et $H_{DR}(V)$ sont donc canoniquement
isomorphes. Etendant les scalaires jusqu'a $\C$, on trouve l'isomorphisme (0.4.1) (le
verifier pour $V = H(X)$).

\subsection*{6.4}
\addcontentsline{toc}{subsection}{6.4}

Soit $E$ une extension finie de $\Q$. La categorie des motifs ``d'Artin a coefficient
dans $E$'' est la categorie deduite de celle des motifs d'Artin comme en 2.1. Nous
l'identifierons a la categorie des $E$-espaces vectoriels de dimension finie, munis
d'une action de $G$.

Les motifs d'Artin a coefficient dans $E$, de rang 1 sur $E$, correspondent aux
caracteres $\chi: G \to E^*$. Nous allons calculer leurs periodes. Soient donc
$\chi: G \to E^*$ et $f$ le conducteur de $\chi$; il se factorise par un caractere, encore note
$\chi$, du quotient $(\Z/f\Z)^* = \Gal(\Q(\exp(2\pi i/f))/\Q)$ de $G$. Notons $[\chi]$ l'espace
vectoriel $E_\chi$, de dimension 1 sur $E$, sur lequel $G$ agit par $\chi$. La somme de Gauss
\[
g = \sum_{u \mod f} \chi(u) \otimes \exp(2\pi i u/f) \in [\chi] \otimes_\Q \overline{\Q}
\]
est non nulle, et invariante par $G$: c'est une base, sur $E$, de $H_{DR}([\chi])$. Le
determinant de $I: H_B([\chi]) \otimes_\Q \C \to H_{DR}([\chi]) \otimes_\Q \C$, calcule dans les bases
$1$ et $g$, vaut $g^{-1}$. On sait que, pour tout plongement complexe $\sigma$ de $E$, on a
$\sigma g \cdot \bar{\sigma} g = f$, nombre rationnel independant de $\sigma$, d'ou
\begin{equation}\label{eq:6.4.1}
\delta([\chi]) \sim \sum_{u \mod f} \chi(u) \otimes \exp(-2\pi i u/f) \in (E \otimes_\Q \C)^*.
\end{equation}

\begin{proposition}\label{prop:6.5}
Soit $D$ le motif $[\chi](n)$. C'est un motif sur $\Q$, a coefficients dans $E$, de rang 1
et de poids $-2n$. Posons $\chi(-1) = (-1)^\eta$, avec $\eta = 0$ ou $1$. On a
\[
\varepsilon^*(D) \sim (2\pi i)^{n + \eta} \cdot \delta(D).
\]
\end{proposition}

On sait que la constante de l'equation fonctionnelle de la fonction $L$ de Dirichlet
$L(\sigma, [\chi], s) = \sum \sigma\chi(n) \cdot n^{-s}$ ($\sigma$ plongement complexe de $E$) est donnee par
\[
\varepsilon(\sigma, [\chi], s) = i^\eta \cdot f^{-s} \cdot \sum_{u \mod f} \sigma\chi(u) \cdot \exp(-2\pi i u/f).
\]
D'apres (5.1.9), on a par ailleurs $\delta(D) = (2\pi i)^n \delta([\chi])$. On conclut en appliquant
(6.4.1) et en notant que pour $s$ entier ($s = n$), $f^{-n}$ est rationnel independant de $\sigma$.

Cette proposition verifie (5.6.1) pour les motifs $[\chi](n)$. Pour achever la preuve de
5.6, il ne reste qu'a enoncer la

\begin{conjecture}\label{conj:6.6}
Tout motif sur $\Q$, a coefficient dans $E$ et de rang 1 est de la forme $[\chi](n)$, pour
un caractere $\chi$ de $G = \Gal(\overline{\Q}/\Q)$ a valeurs dans les racines de l'unite de $E$ et
$n$ un entier.
\end{conjecture}

\begin{proposition}\label{prop:6.7}
La Conjecture 2.8 est vraie pour les fonctions $L$ d'Artin.
\end{proposition}

Pour les motifs d'Artin, on dispose de l'equation fonctionnelle des fonctions $L$.
De plus, le determinant d'un motif d'Artin, tordu a la Tate, est du type predit en
6.6. Les arguments des paragraphes 5, 6 montrent donc la compatibilite de 2.8 a
l'equation fonctionnelle, et il suffit de prouver 2.8 pour les motifs $V(n)$, pour $V$ un
motif d'Artin et $n$ un entier $\geq 0$. Si $V(n)$ est critique, $F_\infty$ agit alors sur $V$ par
multiplication par $-(-1)^n$, et $c^+ = 1$ (cf.~3.1), et il s'agit de prouver que, pour
tout plongement $\sigma$ de $E$ dans $\C$ et tout automorphisme de $G$, on a
$\sigma L(\sigma, V(n)) = L(\sigma, V(n))$.

On le deduit des resultats de Siegel [14]. Voir [2, 1.2].


\section*{7. Fonctions $L$ attachees aux formes modulaires}
\addcontentsline{toc}{section}{7. Fonctions L attachees aux formes modulaires}

\subsection*{7.1}
\addcontentsline{toc}{subsection}{7.1}

Posons $q = e^{2\pi i z}$, et soit $f = \sum a_n q^n$ une forme modulaire holomorphe cuspi-
dale primitive (newform) de poids $k \geq 2$, conducteur $N$ et caractere $\varepsilon$. La serie de
Dirichlet $\sum a_n n^{-s}$ admet un developpement en produit eulerien, de facteur local
en $p \nmid N$ egal a $(1 - a_p p^{-s} + \varepsilon(p) p^{k-1-2s})^{-1}$.

Soit $E$ le sous-corps de $\C$ engendre par les $a_n$. La forme $f$ doit donner lieu a un
motif $M(f)$ de rang 2 a coefficient dans $E$, de type de Hodge $\{(k-1,0), (0,k-1)\}$,
de determinant $[\varepsilon^{-1}](1-k)$ (notations de 6.4) et de fonction $L$ la serie de Dirichlet
$\sum a_n n^{-s}$.

Je n'ai pas essaye de definir $M(f)$ comme etant un motif au sens de Grothendieck.
Une difficulte est que $M(f)$ apparait de facon naturelle comme facteur direct dans
la cohomologie d'une variete non compacte, ou encore comme facteur direct dans
la cohomologie d'une courbe modulaire complete, a coefficient dans l'image directe
d'un faisceau localement constant (plutot, un systeme local de motifs!) sur la partie
a distance finie de cette courbe modulaire. Ceci echappe au formalisme de Gro-
thendieck, mais permet de definir les realisations du motif $M(f)$.

\subsection*{7.2}
\addcontentsline{toc}{subsection}{7.2}

Supposons tout d'abord que $k = 2$, et que $\varepsilon$ est trivial. Soient $X$ le demi-plan
de Poincare, $N$ le conducteur de $f$; $\Gamma_0(N)$ le sous-groupe de $\SL(2, \Z)$ forme
des matrices dont la reduction mod $N$ est de la forme $\begin{pmatrix} * & * \\ 0 & * \end{pmatrix}$,
et posons $\omega_f = \sum a_n q^n \cdot dq/q = \sum a_n q^n \cdot 2\pi i \, dz$. La forme
$\omega_f$ est une forme differentielle holomorphe sur la courbe complete $M(\Gamma_0(N))$ de
$M^0(\Gamma_0(N)) = X/\Gamma_0(N)$. Elle est vecteur propre des correspondances de Hecke:
\begin{equation}\label{eq:7.2.1}
T_n \omega_f = a_n \omega_f \quad \text{(pour $n$ premier a $N$)}
\end{equation}
et est caracterisee a un facteur pres par (7.2.1). Ce fait resulte du theoreme fort de
multiplicite 1 et de la theorie des formes primitives; le lecteur peut, s'il le prefere
completer (7.2.1) par une condition analogue pour $n$ non premier a $N$, et faire de
meme ci-dessous; l'assertion devient alors elementaire, car la condition (7.2.1)
ainsi completee determine (a un facteur pres) le developpement de Taylor de $\omega_f$
en la pointe $i\infty$.

Abrgeons $M^0(\Gamma_0(N))$ en $M^0$ et $M(\Gamma_0(N))$ en $M$. Ces courbes ont une
$\Q$-structure naturelle, pour laquelle les correspondances de Hecke sont definies sur
$\Q$. La forme $\omega_f$ est definie sur $E$, et son conjugue par un automorphisme
$\sigma$ de $\C$ est $\omega_{f^\sigma}$ (appliquer $\sigma$ aux coefficients).

Definissons le motif $M(f)$ comme etant, dans le langage 2.1(B), le sous-motif de
$H^1(M)_E$ noyau des endomorphismes $T_n - a_n$---si on veut ne considerer que des
noyaux de projecteurs, remplacer $T_n - a_n$ par $P(T_n - a_n)^N$ pour $P$ un polynome
convenable. On sait que $M(f)$ a les proprietes enoncees en 7.1.

Dans chacune des theories cohomologiques $\mathcal{H}$ qui nous interessent, l'application
\begin{equation}\label{eq:7.2.2}
\mathcal{H}^0(M^0) \to \mathcal{H}^1(M^0)
\end{equation}
est surjective, et les systemes de valeurs propres des $T_n$ sur le noyau n'apparaissent
pas dans $\mathcal{H}^0(M^0)$. La $\mathcal{H}$-realisation de $M(f)$ est donc encore le noyau commun,
dans $\mathcal{H}^1(M) \otimes_\Q E$, des $T_n - a_n$.

Le groupe de cohomologie $H^1_c(M, \Q)$ est muni d'une structure de Hodge mixte,
de type $\{(0,0), (0,1), (1,0)\}$, dont $H^1(M, \Q)$ est le quotient de type
$\{(0,1), (1,0)\}$. En particulier, (7.2.2) induit un isomorphisme sur les sous-espaces
$F^1$ de la filtration de Hodge; toute forme differentielle holomorphe $\omega$ sur $M$ definit
ainsi une classe de cohomologie a support propre sur $M$. Ceci peut se voir
directement: si $\mathcal{J}$ est l'ideal des pointes, $H_c^*(M^0)$, en cohomologie de de Rham,
est l'hypercohomologie sur $M^{an}$ du complexe $\mathcal{J} \to \Omega^1$ (analytiquement, ce complexe
est une resolution du faisceau constant $\C$ sur $M^0$, prolonge par 0 sur $M$), et le
$H^1$ recoit $H^0(M, \Omega^1)$.

Supposons pour simplifier que $E = \Q$, et calculons $c^+(M(f)(1))$. Le motif $M(f)(1)$
est le dual de $M(f)$ et $c^+$ est donc une periode de $M(f)$ (1.7): il faut integrer
$\omega_f$ contre une classe d'homologie rationnelle de $M$, fixe par $F_\infty$. Relevant
$M(f)$ dans $H^1(M)$, on voit qu'on peut plutot integrer $\omega_f$ contre une classe
d'homologie sans support de $M$, fixe par $F_\infty$. Toutes les integrales non nulles de
ce type seront commensurables.

Pour calculer $F_\infty$, il est utile d'ecrire $M$ comme quotient de
$X = \C - \R$ par le sous-groupe de $\GL(2, \Z)$ forme des matrices ayant pour
reduction mod $N$ une matrice $\begin{pmatrix} * & * \\ 0 & * \end{pmatrix}$. La conjugaison
complexe est alors induite par $z \mapsto \bar{z}$, et l'image dans $M(\C)$ de $i\R^+$ est
un cycle sans support fixe par $F_\infty$. La formule $L(M(f), 1) = \int_{i\R^+} \omega_f$
justifie 1.8 pour $M(f)(1)$: le second membre est, ou bien nul, ou bien $\sim c^+(M(f)(1))$.

\begin{remark}
On a aussi
\[
c^+(M(f)(1)) \sim \int_{i\R^+} \omega_f \sim \int_0^{i\infty} f(z) (iz)^{k-2} \, dz
\]
sous-groupe $\Gamma_1(N)$ de $\SL(2, \Z)$.
\end{remark}

\subsection*{7.4}
\addcontentsline{toc}{subsection}{7.4}

Pour $\varepsilon$ (et $E$) quelconque, il faut remplacer $\Gamma_0(N)$ par $\Gamma_1(N)$: le sous-groupe
de $\SL(2, \Z)$ forme des matrices de reduction mod $N$ de la forme
$\begin{pmatrix} 1 & * \\ 0 & 1 \end{pmatrix}$. Comme ci-dessus, pour calculer $F_\infty$, il est plus commode de
traivailler avec $X = \C - \R$ et le sous-groupe de $\GL(2, \Z)$ forme des matrices
de reduction mod $N$ de la forme $\begin{pmatrix} 1 & * \\ 0 & * \end{pmatrix}$.

Par ailleurs, une dualisation apparait, cachee en 7.2 par la symetrie de la
correspondance $T_n$. Pour une definition convenable de $T_n$, on a
$\langle T_n x, y \rangle = \langle x, T_n y \rangle$ dans $H^1_c(M) \times H_1(M)$ ($M = X/\Gamma_1(N)$), de
sorte que

(a) $M(f)$ est le noyau commun des $T_n - a_n$ dans $H^1(M)_E$;

(b) On a $T_n \omega_f = a_n \omega_f$, et $T_n \omega_f = \varepsilon(n)^{-1} a_n \omega_f$ (noter la formule
$d^\vee = d \cdot \varepsilon^{-1}$), de sorte que $\omega_f$ est dans la realisation de de Rham de $M(f) = M(f)^\vee(-1)$.

On trouve que $c^+(M(f)(1)) \in (E \otimes_\Q \C)^*/E^* \cong \C^{\Hom(E,\C)}/E^*$
est donne par le systeme de periodes
\[
c^+(M(f)(1)) \sim \left(\int_0^{i\infty} f^\sigma(z) (iz)^{k-2} \, dz\right)_\sigma.
\]
Noter que si l'une de ces integrales s'annule, elles s'annulent toutes. Tel est le cas si
et seulement si, dans la partie de l'homologie sans support tendue par le cycle
$i\R^+$ et ses transformes par les $T_n$, le systeme de valeurs propres $a_n$ pour les
$T_n$, $n$ premier a $N$, n'apparait pas. Ceci justifie 2.3, et, partiellement 2.7, pour
$M(f)(1)$.

\begin{remark}
Les systemes de valeurs propres des $T_n$ dans
$\Ker(H^0_c(M) \to H^0(M))$ sont lies aux series d'Eisenstein, alors que ceux qui
apparaissent dans $H^0_c(M)$ sont lies aux formes paraboliques. C'est pourquoi ces
ensembles sont disjoints. Il en resulte que la structure de Hodge mixte de
$H^1_c(M, \Q)$ est somme de structures de Hodge: l'extension de $H^1(M, \Q)$ par
$\Ker(H^1_c(M, \Q) \to H^1(M, \Q))$ splitte. Ceci, generalise au cas de n'importe
quel sous-groupe de congruence de $\SL(2, \Z)$, equivaut au theoreme de Manin [9]
selon lequel la difference entre deux pointes est toujours d'ordre fini dans la
jacobienne (cf.~[6, 10.3.4, 10.3.8 et 10.1.3]). La demonstration donnee n'en differe
d'ailleurs pas en substance de celle de Manin.
\end{remark}

\subsection*{7.6}
\addcontentsline{toc}{subsection}{7.6}

En poids $k$ quelconque, il faut remplacer la cohomologie de $M$ par la cohomologie
de $M$ a coefficient dans un faisceau convenable. Pour decrire ce qui se passe, je
remplacerai $M^0$ et $M$ par $M^0_n$ et $M_n$, relatifs au groupe de congruence
$\Gamma(n)$, avec $n$ multiple de $N$ et $\geq 3$. On redescend ensuite a $M^0$ et $M$ en prenant
les invariants par un groupe fini convenable. Soient $g: E \to M^0_n$ la courbe ellipti-
que universelle et $j$ l'inclusion de $M^0_n$ dans $M_n$. La cohomologie a considerer est
$H^1(M_n, j_* \Sym^{k-2}(R^1 g_* \Q))$. Les realisations de $M(f)$ sont facteur direct de ce
groupe, calcule dans la theorie de cohomologie correspondante. On peut comme
precedemment les relever dans $H^1_c(M^0_n, \Sym^{k-2}(R^1 g_* \Q))$. Pour calculer $c^+$ du
dual de $M(f)$, il faut alors integrer la classe definie par $f$ contre des classes
d'homologie sans support de $M^0_n$ a coefficient dans le systeme local dual de
$\Sym^{k-2}(R^1 g_* \Q)$. Si on prend le cycle image de $i\R^+$, muni de diverses sections
du dual (une base), on trouve les integrales d'Eichler
\[
r_{f,l} = \int_0^{i\infty} f(z) (iz)^l \, dz \quad (0 \leq l \leq k-2)
\]
(periodes paires pour $l$ pair, impaires pour $l$ impair)---et on sait comment ecrire
$L(M(f), n)$ pour $n$ critique en terme de ces periodes.

\begin{proposition}\label{prop:7.7}
Soit $M$ un motif de rang 2, a coefficient dans $E$, de type de Hodge
$\{(a,b), (b,a)\}$ avec $a \neq b$. Alors, $d^\pm(\Sym^n M)$ et $c^\pm(\Sym^n M)$ sont donnes
par les formules suivantes:

$(1)$ si $n = 2l+1$: $d^+ = l+1$ et
\begin{align*}
c^+(\Sym^n M) &= c^+(M)^{l+1} \cdot c^-(M)^l \cdot \delta(M)^{l(l+1)/2},\\
c^-(\Sym^n M) &= c^+(M)^l \cdot c^-(M)^{l+1} \cdot \delta(M)^{l(l+1)/2}.
\end{align*}

$(2)$ si $n = 2l$: $d^+ = l+1$, $d^- = l$ et
\begin{align*}
c^+(\Sym^n M) &= (c^+(M) c^-(M))^{l(l+1)/2} \cdot \delta(M)^{l(l+1)/2},\\
c^-(\Sym^n M) &= (c^+(M) c^-(M))^{l(l+1)/2} \cdot \delta(M)^{l(l-1)/2}.
\end{align*}
\end{proposition}

C'est une question de simple algebre lineaire, que nous traiterons par ``analyse
dimensionnelle''.

(a) $H_B(M)$ est un espace vectoriel de rang 2 sur $E$, muni d'une involution $F_\infty$,
de la forme $\begin{pmatrix} 1 & 0 \\ 0 & -1 \end{pmatrix}$ dans une base $e_+, e_-$ convenable. Les
$d^\pm(\Sym^n M)$ sont les dimensions des parties $+$ et $-$ de la puissance symetrique
$n$-ieme, de bases respectives
$\{e_+^n, e_+^{n-1}e_-, \ldots\}$ et $\{e_+^{n-1}e_-, e_+^{n-3}e_-^3, \ldots\}$---d'ou les valeurs annoncees.

(b) Soit $\omega, \eta$ une base du dual de $H_{DR}(M)$, telle que $\omega$ annule $F^+ H_{DR}(M)$. Le
sous-espace $F^{n \cdot a}$ de $H_{DR}(\Sym^n M)^\vee \cong \Sym^n H_{DR}(M)^\vee$ admet pour base
les $\omega^n, \omega^{n-1}\eta, \ldots, \omega^{n-a+1}\eta^{a-1}$. Et
\begin{align*}
c^+(\Sym^n M) &= \langle \omega^n \wedge (\omega^{n-1}\eta) \wedge \cdots \rangle_{e_+^n, e_+^{n-1}e_-, \ldots} \\
c^-(\Sym^n M) &= \langle \omega^n \wedge (\omega^{n-1}\eta) \wedge \cdots \rangle_{e_+^{n-1}e_-, e_+^{n-3}e_-^3, \ldots}. \tag{7.7.1} \label{eq:7.7.1}
\end{align*}

(c) Posons $V = H_B(M) \otimes_\Q \C \cong H_{DR}(M) \otimes_\Q \C$; c'est un
$E \otimes_\Q \C$-module. Les formules ci-dessus montrent que $c^\pm(\Sym^n M)$ ne
depend que du $E \otimes_\Q \C$-module $V$, de sa base $e_+, e_-$, de $\omega \in V^*$, et de
l'image $\bar{\eta}$ de $\eta$ dans $V^*/\langle \omega \rangle$: le $di$-vecteur a gauche du produit
scalaire (7.7.1) ne change pas si on remplace $\eta$ par $\eta + \lambda\omega$. Par ailleurs,
$\langle \omega, e_+ \rangle$ et $\langle \omega, e_- \rangle$ sont inversibles, et $\bar{\eta}$ est une base de
$V^*/\langle \omega \rangle$. Le systeme $(V, e_+, e_-, \omega, \bar{\eta})$ est donc decrit a isomorphisme pres
par les quantites $c^+ = \langle \omega, e_+ \rangle$, $c^- = \langle \omega, e_- \rangle$ et
$\delta = \langle \omega \wedge \eta, e_+ \wedge e_- \rangle$, dans $(E \otimes \C)^*$. Remplacer $e_+, e_-,
\omega, \eta$ par $\lambda e_+, \mu e_-, \alpha \omega, \nu \eta/\lambda\mu$ remplace $c^+, c^-$ et $\delta$ par
$\lambda\alpha c^+$, $\mu\alpha c^-$ et $\nu\delta$. Pour $c^+ = c^- = \delta = 1$, le second membre de
(7.7.1) est dans $\Q^*$. Dans le cas general $c^\pm(\Sym^n M)$ est donc un multiple
rationnel du produit de $c^+(M)$, $c^-(M)$ et $\delta(M)/c^+(M)c^-(M)$ aux puissances
respectives les degres auxquels figurent $e_+, e_-$ et $\bar{\eta}$ dans (7.7.1). On s'epargne
la moitie du calcul en notant que remplacer $F_\infty$ par $-F_\infty$ echange $e_+$ et
$e_-$, donc $c^+$ et $c^-$, respecte $\delta$, et echange les (7.7.1)$^\pm$ pour $n$
impair, les conserve pour $n$ pair.

$n = 2l+1$: $d^+ = d^- = l+1$,
\begin{align*}
\deg \bar{\eta} \text{ dans } (7.7.1)^+ &= 0 + 1 + \cdots + l = l(l+1)/2,\\
\deg e_+^{n-2i} \text{ dans } (7.7.1)^+ &= (2l+1) + (2l-1) + \cdots + 1 = (l+1)^2\\
&= \deg e_-^{2i+1} \text{ dans } (7.7.1)^-;\\
\deg e_+^{n-2i-1} \text{ dans } (7.7.1)^- &= 2l + (2l-2) + \cdots + 0 = l(l+1)\\
&= \deg e_-^{2i} \text{ dans } (7.7.1)^+.
\end{align*}

$n = 2l$: $d^+ = l+1$, $d^- = l$,
\begin{align*}
\deg \bar{\eta} \text{ dans } (7.7.1)^+ &= 0 + 1 + \cdots + l = l(l+1)/2,\\
\deg \bar{\eta} \text{ dans } (7.7.1)^- &= 0 + 1 + \cdots + (l-1) = l(l-1)/2,\\
\deg e_+ \text{ dans } (7.7.1)^+ &= 2l + (2l-2) + \cdots + 0 = l(l+1),\\
\deg e_+ \text{ dans } (7.7.1)^- &= (2l-1) + (2l-3) + \cdots + 1 = l^2.
\end{align*}

\subsection*{7.8}
\addcontentsline{toc}{subsection}{7.8}

Cette proposition nous fournit une conjecture pour les valeurs aux entiers
critiques de $L(\Sym^n M(f), s)$. Voici les formules:
\begin{equation}\label{eq:7.8.1}
L(\sigma, M(f), s) = \sum a_n^\sigma n^{-s} = \prod_p L_p(\sigma, M(f), s)
\end{equation}
o\`u pour presque tout $p$
\[
L_p(\sigma, M(f), s)^{-1} = (1 - a_p^\sigma p^{-s})(1 - \bar{a}_p^\sigma p^{-s})
\]
avec $\bar{\sigma}$ un caractere de Dirichlet. On a $\bigwedge^2 M(f) = [\varepsilon^{-1}](1-k)$, d'ou
\begin{equation}\label{eq:7.8.2}
\delta(M(f)) \sim (2\pi i)^{1-k} \varepsilon(f) \otimes \exp(-2\pi i/p),
\end{equation}
si $\varepsilon$ est de conducteur $F$, et
\begin{equation}\label{eq:7.8.3}
L^*(M(f), m) \sim (2\pi i)^m c^+(M(f)), \quad m = (-1)^m, \quad \text{pour } 1 \leq m \leq k-1.
\end{equation}
\begin{equation}\label{eq:7.8.4}
L(\sigma, \Sym^n M(f), s) = \prod_p L_p(\sigma, \Sym^n M(f), s)
\end{equation}
o\`u pour presque tout $p$
\[
L_p(\Sym^n M(f), s)^{-1} = \prod_{i=0}^n (1 - \alpha_p^i \beta_p^{n-i} p^{-s});
\]
conjecturalement, pour $m$ critique, on a
\[
L^*(\Sym^n M(f), m) \sim (2\pi i)^{m \cdot d^+(\Sym^n M(f))} \cdot c^+(\Sym^n M(f)),
\]
o\`u $d^+$ et $c^+$ sont donnes en terme de $c^\pm(M(f))$ et $\delta(M(f))$ (caracterises
par (7.8.2), (7.8.3)) par les formules 7.7.

Pour l'evidence numerique en faveur de cette conjecture, voir [18].


\section*{8. Caracteres de Hecke algebriques}
\addcontentsline{toc}{section}{8. Caracteres de Hecke algebriques}

Le lecteur trouvera dans [3, paragraphe 5], dont nous utiliserons les notations, les
d\'efinitions essentielles relatives aux caract\`eres de Hecke alg\'ebriques (= Gr\"ossencharaktere
 de type $A_0$).

\begin{conjecture}\label{conj:8.1}
Soient $k$ et $E$ deux extensions finies de $\Q$, et $\chi$ un caract\`ere de Hecke alg\'ebrique
 de $k$ \`a valeurs dans $E$.

\textup{(i)} Il existe un motif $M(\chi)$ de rang $[k:\Q]$ sur $k$, \`a coefficients dans $E$, tel que,
pour toute place $\lambda$ de $E$, la repr\'esentation $\lambda$-adique $H_\lambda M(\chi)$ soit celle
d\'efinie par $\chi$: le Frobenius g\'eom\'etrique en $\mathcal{P}$ premier au conducteur de $\chi$ et \`a la
caract\'eristique r\'esiduelle $l$ de $\lambda$ agit par multiplication par $\chi(\mathcal{P})$.

\textup{(ii)} Ce motif est caract\'eris\'e \`a isomorphisme pr\`es par cette propri\'et\'e.

\textup{(iii)} Tout motif de rang 1 est de la forme $M(\chi)$.

\textup{(iv)} D\'ecomposons $k \otimes_\Q E$ en produit de corps: $k \otimes_\Q E = \prod K_i$, et \'ecrivons la partie
alg\'ebrique $\chi_{\text{alg}}: k^* \to E^*$ de $\chi$; sous la forme
$\chi_{\text{alg}}(x) = \prod N_{K_i/E}(x)^{n_i}$. La d\'ecomposition de $k \otimes_\Q E$ induit une d\'ecomposition
de $H_{DR}(M(\chi))$ en les $H_{DR}(M(\chi))_i = H_{DR}(M(\chi)) \otimes_{k \otimes E} K_i$;
avec cette notation, la filtration de Hodge est la filtration par les degr\'es:
$F^p H_{DR}(M(\chi)) = \bigoplus_{n_i \geq p} H_{DR}(M(\chi))_i$.
\end{conjecture}


L'unicit\'e 8.1(ii) impose aux $M(\chi)$ le formalisme suivant
\begin{equation}\label{eq:8.1.1}
M(\chi' \chi'') \sim M(\chi') \otimes M(\chi'').
\end{equation}
(8.1.2) Si $\xi: E \hookrightarrow E'$ est une extension finie de $E$, $M(\xi\chi)$ se d\'eduit de $M(\chi)$
par extension des coefficients de $E$ \`a $E'$.

(8.1.3) Si $k'$ est une extension finie de $k$, $M(\chi \circ N_{k'/k})$ se d\'eduit de $M(\chi)$
par extension des scalaires de $k$ \`a $k'$.

\begin{remark}
Posons les notations: $c =$ la conjugaison complexe, $\overline{\Q} =$ la cl\^oture alg\'ebrique de
$\Q$ dans $\C$, $S = \Hom(k, \overline{\Q}) = \Hom(k, \C)$, $J = \Hom(E, \overline{\Q}) = \Hom(E, \C)$. La
d\'ecomposition de $k \otimes_\Q E$ en les $K_i$-corps correspond \`a la partition de
$S \times J$ en les orbites de $\Gal(\overline{\Q}/\Q)$.

Tout homomorphisme alg\'ebrique $\eta: k^* \to E^*$ s'\'ecrit sous la forme
$\prod N_{K_i/E}(x)^{n_i}$. Si l'orbite sous $\Gal(\overline{\Q}/\Q)$ de $(\sigma, \tau) \in S \times J$
correspond \`a $K_i$, i.e., si $\sigma \otimes \tau: k \otimes_\Q E \to \C$ se factorise par $K_i$,
nous noterons $n(\eta; \sigma, \tau)$, ou simplement $n(\sigma, \tau)$, l'entier $n_i$. La fonction
$n(\sigma, \tau)$ est constante sur les orbites de $\Gal(\overline{\Q}/\Q)$. Si $\eta$ est la partie
alg\'ebrique d'un caract\`ere de Hecke $\chi$: on a de plus
\begin{equation}\label{eq:8.2.1}
\text{L'entier } w = n(\eta; c\sigma, \tau) + n(\eta; \sigma, \tau) \text{ est ind\'ependant de } \sigma \text{ et } \tau.
\end{equation}

On \'ecrira parfois $n(\chi; \sigma, \tau)$ pour $n(\chi_{\text{alg}}; \sigma, \tau)$.

C'est le poids de $\chi$ (et de $M(\chi)$).

R\'eciproquement, un homomorphisme $\eta$ v\'erifiant (8.2.1) est presque de la forme
$\chi_{\text{alg}}$, pour $\chi$ un caract\`ere de Hecke convenable:

(a) une de ses puissances l'est;

(b) il existe une extension finie $\xi: E \hookrightarrow E'$ de $E$ telle que $\xi \circ \eta$ le soit;

(c) il existe une extension finie $k'$ de $k$ telle que $\eta \circ N_{k'/k}$ le soit.

La r\`egle 8.1(iv) permet de d\'eduire la bigraduation de Hodge de $H_\sigma M(\chi)$ de sa
structure de $E$-module: le facteur direct $H_\sigma M(\chi) \otimes_{E,\tau} \C$ de
$H_\sigma M(\chi) \otimes_\Q \C$ est de type de Hodge $(p,q)$, avec
$p = n(\chi; \sigma, \tau)$ et $q = n(\chi; c\sigma, \tau) = w - p$.
\end{remark}


\begin{example}
Soit $A$ une variete abelienne sur $k$, a multiplication complexe par $E$. On suppose
$H^1(A)$ de rang 1 sur $E$ (type CM). Il resulte de la theorie de Shimura et
Tanayama que $H_\sigma(A)$ verifie la condition de 8.1(i), pour un caractere de Hecke
algebrique $\chi$ de $k$ a valeurs dans $E$, et que la partie algebrique $\chi_{\text{alg}}$ de $\chi$
agit sur le $k \otimes_\Q E$-module $\text{Lie}(A)$:
\[
\chi_{\text{alg}}(x) = \det_{k \otimes E}(x \otimes 1, \text{Lie}(A)).
\]
\end{example}

\begin{example}
Prenons pour $k$ le corps des racines $n$-iemes de l'unite, soit $V$ l'hypersurface
de Fermat d'equation projective $\sum_{i=0}^{m} X_i^n = 0$ et soit $M$ le motif ``cohomologie
primitive de dimension moitie de $V$''. Soit $G$ le quotient de $\mu_n^{m+1}$ par son
sous-groupe diagonal. Ce groupe agit sur $V$ par $(\alpha_i) \cdot (X_i) = (\alpha_i X_i)$, et
sur $M$ par transport de structure. Decomposant $M$ a l'aide de la decomposition de
l'algebre de groupe $\Q[G]$ en produit de corps, on obtient des motifs verifiant 8.1(i)
pour des caracteres de Hecke algebriques convenables: ceux introduits par Weil dans
son etude des sommes de Jacobi.
\end{example}

\begin{example}
Le motif $\Z(-1)$ verifie 8.1(i) pour $\chi =$ la norme. Pour $\chi$ d'ordre fini, un motif
d'Artin convenable verifie 8.1(i).
\end{example}

\subsection*{8.6}
\addcontentsline{toc}{subsection}{8.6}

Pour chaque $\sigma \in S$, $H_\sigma(M(\chi))$ est de rang 1 sur $E$. Choisissons une base $e_\sigma$
de chaque $H_\sigma(M(\chi))$. Outre sa structure de $E$-module, la somme des
$H_\sigma(M(\chi)) \otimes \C$ a une structure naturelle de $k \otimes_\Q E \otimes_\Q \C = \C^{S \times J}$-module
libre de rang 1, pour laquelle $e = \sum_\sigma e_\sigma$ est une base.

\subsection*{8.7}
\addcontentsline{toc}{subsection}{8.7}

La realisation de de Rham $H_{DR}(M(\chi))$ est un $k \otimes_\Q E$-module libre de rang 1.
Choisissons en une base $\omega$. La somme des $I_\sigma: H_\sigma(M(\chi)) \otimes_\Q \C \to
H_{DR}(M(\chi)) \otimes_\Q \C$ est l'isomorphisme (0.4) de $k \otimes_\Q E \otimes_\Q \C$-modules
\[
\bigoplus_\sigma H_\sigma(M(\chi)) \otimes_\Q \C \xrightarrow{\sim} H_{DR}(M(\chi)) \otimes_\Q \C.
\]

Soit $\omega$ une base du $k \otimes_\Q E$-module $H_{DR}(M(\chi))$, et posons
$p'(\chi) = \omega/I(e) \in (k \otimes_\Q E \otimes_\Q \C)^*$. Cette periode depend de $\chi$, $\omega$ et $e$.
Prise modulo $(E \otimes_\Q F)^*$ et $E^*$, elle ne depend que de $\chi$. Nous noterons
$p'(\chi; \sigma, \tau)$---ou simplement $p'(\sigma, \tau)$---la composante d'indice $(\sigma, \tau)$ de
son image par l'isomorphisme $E \otimes_\Q F \otimes_\Q \C \xrightarrow{\sim} \C^{S \times J}$.

\subsection*{8.8}
\addcontentsline{toc}{subsection}{8.8}

Soit $A$ comme en 8.3, et calculons les periodes $p'(\sigma, \tau)$, modulo $\Q^*$, pour le
motif $H^1(A)$. On commence par etendre les scalaires de $k$ a $\overline{\Q}$, a l'aide de $\sigma$.
Si $\tau(1) = 1$, il existe alors une 1-forme holomorphe $\omega$ definie sur $\overline{\Q}$ telle que
$u^* \omega = \tau(u) \omega$ pour $u \in E$, et, pour $z \in H_1(A(\C))$, on a $p'(\sigma, \tau) \sim 2\pi i$.
On peut passer de la au cas general a l'aide de la formule
$p'(\sigma, \tau) \cdot p'(c\sigma, \tau) \sim 2\pi i$.

\subsection*{8.9}
\addcontentsline{toc}{subsection}{8.9}

La conjecture 8.1(ii) affirme en particulier que si deux motifs verifient la
condition de 8.1(i), ils ont meme periode $p'$. Pour les motifs 8.4, les periodes
s'expriment en terme de valeurs de la fonction $\Gamma$ et, si on travaille mod $\Q^*$,
8.1 suggere la conjecture de B.~Gross [7] reliant certaines periodes a des produits
de valeurs de la fonction $\Gamma$.

La comparaison de 8.4 et 8.5 mene a la Conjecture 8.11, 8.13 suivante. Le resultat
annonce apres 0.10 permet de la demontrer.

\subsection*{8.10}
\addcontentsline{toc}{subsection}{8.10}

Soient $N$ un entier, $k = \Q(\exp(2\pi i/N))$, $\mathcal{P}$ un ideal premier de $k$, premier
a $N$, $k_\mathcal{P}$ le corps residuel et $q = N\mathcal{P} = |k_\mathcal{P}|$. On notera $t$ l'inverse de la reduction
mod $\mathcal{P}$, des racines $N$-iemes de 1 dans $k_\mathcal{P}$ a celles de $k$.

Pour $a \in N^{-1}\Z/\Z$, $0 \leq a \leq 1$, considerons la somme de Gauss
$g(\mathcal{P}, a, \psi) = -\sum_{x \in k_\mathcal{P}^*} t(x^{-a(q-1)}) \psi(\Tr(x))$. La somme est etendue a
$k_\mathcal{P}^*$ et $\psi: k_\mathcal{P} \to \C^*$ est un caractere additif non trivial.

Soit $a = \sum n(a) \delta_a$ dans le groupe abelien libre de base $N^{-1}\Z/\Z - \{0\}$. Si
$\sum n(a) \cdot a = 0$, le produit des $g(\mathcal{P}, a, \psi^{u(a)})$ est independant de $\psi$ et on pose
$g(\mathcal{P}, a) = \prod g(\mathcal{P}, a, \psi^{u(a)})$. Weil [16] a montre que, comme fonction de $\mathcal{P}$,
$g(\mathcal{P}, a)$ est un caractere de Hecke algebrique $\chi_a$ de $k$ a valeurs dans $k$. Notons
$(a)$ le representant entre 0 et 1 de $a$ dans $\Z/N$. Si $a = \sum n(a) \delta_a$ verifie
\begin{equation}\label{eq:star}
\text{Pour tout } u \in (\Z/N)^*\text{, on a } \sum n(a) \cdot (ua) = 0,\tag{$*$}
\end{equation}
il resulte de la determination par Weil de la partie algebrique de $\chi_a$ que
$\chi_a$ est d'ordre fini; on note encore $\chi_a$ le caractere de $\Gal(\overline{\Q}/k)$ valant
$\chi_a(\mathcal{P})$ sur le Frobenius geometrique en $\mathcal{P}$.

Posons $\Gamma(a) = \prod \Gamma((a))^{n(a)}$. Une premiere partie de la conjecture:
algebricite de $\Gamma(a)$ lorsque $a$ verifie $(*)$, a ete prouvee par Koblitz et Ogus:
voir l'appendice. On desire avoir, plus precisement:

\begin{conjecture}\label{conj:8.11}
Si $a$ verifie $(*)$, on a $\sigma \Gamma(a) = \chi_a(\sigma)^{-1} \Gamma(a)$ pour tout
$\sigma \in \Gal(\overline{\Q}/k)$.
\end{conjecture}

Si $a$ est invariant par un sous-groupe $H$ de $(\Z/N)^*$, on peut preciser 8.11 en
remplacant $\Gal(\overline{\Q}/k)$ par $\Gal(\overline{\Q}/k^H)$, et en utilisant SGA 4$^{1/2}$ 6.5 pour definir un
caractere de ce groupe, a valeur dans les racines de l'unite de $k^*$. C'est ce que nous
faisons ci-dessous.

\subsection*{8.12}
\addcontentsline{toc}{subsection}{8.12}

Soient $H$ un sous-groupe de $(\Z/N)^*$, $\mathcal{P}$ un ideal premier de $k^H$, premier a $N$,
$\mathbb{F}$ son corps residuel et $\psi: \mathbb{F} \to \C^*$ un caractere additif non trivial. Pour
$a \in N^{-1}\Z/\Z - \{0\}$, soient $\mathcal{P}_\alpha$ les ideaux premiers de $k^H(\exp(2\pi i a))$
au-dessus de $\mathcal{P}$. Si $\mathbb{F}'$ est le corps residuel en $\mathcal{P}_\alpha$, et que $[\mathbb{F}':\mathbb{F}] = (1)$, on
note $g(\mathcal{P}_\alpha, a, \psi)$ la somme de Gauss $-\sum t(x^{-a(q'-1)}) \psi(\Tr_{\mathbb{F}'/\mathbb{F}} x)$
(somme etendue a $\mathbb{F}'^*$). Le produit $g(\mathcal{P}, a, \psi) = \prod_\alpha g(\mathcal{P}_\alpha, a, \psi)$ ne
depend que de l'orbite de $a$ sous $H$.

Soit $a = \sum n(a) \delta_a$ invariant par $H$, et verifiant $\sum n(a) \cdot a = 0$. Pour chaque
orbite $\mathcal{O}$ de $H$ dans $N^{-1}\Z/\Z - \{0\}$, on note $n(\mathcal{O})$ la valeur constante de $n(a)$ sur
$\mathcal{O}$. On pose
\[
g(\mathcal{P}, a) = \prod_{\mathcal{O}} g(\mathcal{P}, a, \psi)^{n(\mathcal{O})}.
\]

Comme fonction de $\mathcal{P}$, $g(\mathcal{P}, a)$ est un caractere de Hecke algebrique $\chi_a$ de $k^H$ a
valeurs dans $k^*$. Pour le prouver, on se ramene par additivite a supposer les
$n(a) \geq 0$. On applique alors [3, 6.5] pour $F = k^H$, $F' = \Q$, $k =$ notre $k$, et
$I =$ une somme disjointes de copies d'orbites $H$: $n(\mathcal{O})$ copies de $\mathcal{O}$; pour $i \in I$,
d'image $a/N$ dans $N^{-1}\Z/\Z$, on prend pour $\chi_i$ le compose
$\Z(1) \to (\Z/N)(1) = \mu_N(k) \hookrightarrow k^*$; le caractere de Hecke obtenu est le produit
de $\chi_a$ ci-dessus par le caractere ``signature de la representation de permutation
de $H$ sur $I$''. Un cas particulier de ce resultat figure deja dans Weil [17]. Le
caractere $\chi_a$ de 8.10 est le compose de $\chi_a$ ci-dessus avec la norme $N_{k^H/\Q}$.

\begin{conjecture}\label{conj:8.13}
Soient $\mathcal{P}$ un ideal premier de $k^H$, premier a $N$ et $\mathcal{F}$ le Frobenius
geometrique en $\mathcal{P}$. Si $a$ invariant sous $H$ verifie $(*)$, on a
\[
\mathcal{F} \Gamma(a) = \chi_a(\mathcal{P})^{-1} \Gamma(a) \cdot \prod_{\mathcal{O}} n(\mathcal{O})^{n(\mathcal{O})}.
\]
En particulier, si le groupe des racines de l'unite de $k^H$ est d'ordre $N'$, on a
$\Gamma(a)^{N'} \in k^*$.
\end{conjecture}

On peut de 8.13 deduire la variante suivante, apparemment plus generale. On
remplace la condition $(*)$ par
\begin{equation}\label{eq:starprime}
\sum n(a) \cdot (ua) = k \text{ est un entier independant de } u \in (\Z/N)^*.\tag{$*'$}
\end{equation}
Le caractere de Hecke $N\mathcal{P}^{-k} \cdot g(\mathcal{P}, a)$ est alors d'ordre fini. On l'identifie a un
caractere $\chi$ de $\Gal(\overline{\Q}/k^H)$, et on espere avoir
\[
\sigma((2\pi i)^{-k} \Gamma(a)) = \chi(\sigma)^{-1} ((2\pi i)^{-k} \Gamma(a)).
\]

\subsection*{8.14}
\addcontentsline{toc}{subsection}{8.14}

Si $k$ est un corps extension quadratique totalement imaginaire d'un corps
totalement reel, soit, comme nous dirons, un corps de type CM, Shimura [13] a
determine les valeurs critiques des fonctions $L$ de caracteres de Hecke algebriques
 de $k$, au produit pres par un nombre algebrique. Il les exprime en terme de perio-
des de varietes abeliennes de type CM, a multiplication complexe par $k$. Dans la
fin du paragraphe, nous montrons que son theoreme est compatible a la Conjecture
2.8, qui les exprime en terme de periodes de motifs sur $k$, de rang 1.

\subsection*{8.15}
\addcontentsline{toc}{subsection}{8.15}

Soient $\chi$ et $M(\chi)$ comme en 8.1. Excluons le cas 8.5, et supposons que
$R_{k/\Q} M(\chi)$ verifie 1.7. Le corps $k$ est alors totalement imaginaire, et les nombres
de Hodge $h^{p,p}$ sont nuls: avec les notations de 8.1 et 3.2, aucun $n_i$ n'est egal a $w/2$.
Notre premier tache est de calculer $c^+ R_{k/\Q} M(\chi)$ en terme des periodes
$p'(\chi; \sigma, \tau)$.

Rappelons que c'est le determinant de
\[
I^+: H_B^+ R_{k/\Q} M(\chi) \otimes_\Q \C \xrightarrow{\sim} H_{DR}^+ R_{k/\Q} M(\chi) \otimes_\Q \C,
\]
calcule dans des bases definies sur $E$.

Choisissons les $e_\sigma$ de 8.6 de sorte que $F_\infty e_\sigma I = e_{c\sigma}$. Les $e_\sigma$
($\sigma \in S$), i.e.~$e_\sigma + e_{c\sigma}$, forment alors une base de $H_B^+ R_{k/\Q} M(\chi) \subset
H_B R_{k/\Q} M(\chi) = \bigoplus_{\sigma \in S} H_\sigma M(\chi)$. Notons au passage que
$d^+ = d^- = \tfrac{1}{2}[k:\Q]$. Soit $\bar{S}$ le quotient de $S$ par $\Gal(\C/\R)$. Pour calculer
$c^+ = \det(I^+)$, nous utiliserons la base $(e_\sigma + e_{c\sigma})$ de $H_B^+$. Elle est indexee
par $\bar{S}$.

D'apres 8.1(iv), la filtration de Hodge de $H_{DR} R_{k/\Q} M(\chi) = H_{DR} M(\chi)$ se lit
sur sa structure de $k \otimes_\Q E$-module: si on note $(k \otimes_\Q E)_i$ le facteur direct de
$k \otimes_\Q E$ produit des $K_i$ tels que $n_i < w/2$, le quotient $H_{DR}^+ R_{k/\Q} M(\chi)$
de $H_{DR} R_{k/\Q} M(\chi)$ est le facteur direct correspondant:
\[
H_{DR}^+ R_{k/\Q} M(\chi) = H_{DR} M(\chi) \otimes_{k \otimes E} (k \otimes E)_+.
\]

Soit $\omega$ comme en 8.7, et utilisons la structure de
$k \otimes_\Q E \otimes_\Q \C = \C^{S \times J}$-module de $H_{DR} M(\chi)_\C$ pour decomposer
$\omega$: $\omega = \sum \omega_{\sigma,\tau}$. On a par definition $I(e_\sigma) =
p'(\chi; \sigma, \tau)^{-1} \omega_{\sigma,\tau}$ et donc
\[
I^+(e_\sigma + e_{c\sigma}) = \sum_{\substack{\tau \\ n(\sigma,\tau) < w/2}} p'(\sigma, \tau)^{-1} \omega_{\sigma,\tau}
+ \sum_{\substack{\tau \\ n(c\sigma,\tau) < w/2}} p'(c\sigma, \tau)^{-1} \omega_{c\sigma,\tau};
\]
c'est la somme, indexee par $\tau \in J$, de termes egaux a $p'(\sigma, \tau)^{-1} \omega_{\sigma,\tau}$
pour $n(\sigma, \tau) < w/2$, et a $p'(c\sigma, \tau)^{-1} \omega_{\sigma,\tau}$ pour
$n(c\sigma, \tau) < w/2$ (i.e.~$n(\sigma, \tau) > w/2$). Pour $\sigma \in S$, de representant
$\bar{\sigma}$, posons
\[
\omega_{\bar{\sigma}} = \sum_{\substack{\tau \\ n(\sigma,\tau) < w/2}} \omega_{\sigma,\tau}
+ \sum_{\substack{\tau \\ n(c\sigma,\tau) < w/2}} \omega_{c\sigma,\tau}.
\]
Des $I^+(e_\sigma + e_{c\sigma})$ sont des multiples des $\omega_{\bar{\sigma}}$ par des elements de
$E \otimes_\C$; les $\omega_{\bar{\sigma}}$ forment donc une base de $H_{DR}^+$. Dans les bases
$e_\sigma + e_{c\sigma}$ et $\omega_{\bar{\sigma}}$, la matrice de $I^+$ est diagonale; son determinant
$\det'(I^+) \in (E \otimes_\Q \C)^* = \C^{\Hom(E,\C)}$ a pour coordonnees
\[
\det'(I^+)_\tau = \prod_{n(\sigma,\tau) < w/2} p'(\sigma, \tau)^{-1}.
\]

Soient les applications $E^{\bar{S}} \to E^{\bar{S}}$: $I_\tau \mapsto I_\tau + I_{c\tau}$, pour $\sigma$ image de
$\bar{\sigma}$, $E^{\bar{S}} \otimes_\Q \C \to (k \otimes_\Q E)_+ \otimes_\Q \C$, deduit de l'isomorphisme de
$k \otimes_\Q \C$ avec $\C^S$, et la projection de $k \otimes_\Q E$ sur $(k \otimes_\Q E)_+$. Par
composition, on obtient un isomorphisme de $E \otimes_\Q \C$-modules:
$E^{\bar{S}} \otimes_\Q \C \xrightarrow{\sim} (k \otimes_\Q E)_+ \otimes_\Q \C$.

Nous noterons $D(\chi)$ son determinant, calcule dans des bases definies sur $E$ des
deux membres. Identifiant $H_{DR}^+$ a $(k \otimes_\Q E)_+$, a l'aide de la base $\omega$,
on voit que c'est le determinant de l'application identique de $H_{DR}^+ \otimes_\Q \C$, calcule
dans la base $\omega_{\bar{\sigma}}$ a la source, et une base definie sur $E$, au but. Notant
$\partial_\tau$ ses composantes dans $\C^J$, on trouve pour $c^+ = \det'(I^+) \cdot D(\chi)$ la
formule suivante.

\begin{proposition}\label{prop:8.16}
On a
\[
c^+ R_{k/\Q} M(\chi) = \left(\prod_{n(\sigma,\tau) < w/2} p'(\chi; \sigma, \tau)^{-1} \cdot D(\chi)\right)_{\tau \in J}.
\]
\end{proposition}

\begin{remark}
Supposons $k$ de type CM, extension quadratique de $k_0$ totalement reel. Le quotient
$\bar{S}$ de $S$ s'identifie alors a l'ensemble des plongements complexes de $k_0$, et le
diagramme
\[
\begin{CD}
\bar{S} @>{\sigma \mapsto (\sigma,\tau)}>> k_0 \otimes_\Q E \otimes_\C \\
@VVV @VVV \\
E^{\bar{S}} \otimes_\Q \C @>{\sim}>> (k \otimes_\Q E)_+ \otimes_\Q \C
\end{CD}
\]
est commutatif. L'application composee $k_0 \otimes_\Q E \to (k \otimes_\Q E)_+$ est donc
un isomorphisme, et $D(\chi)$ est encore le determinant de $E^{\bar{S}} \otimes_\Q \C \xrightarrow{\sim}
k_0 \otimes_\Q E \otimes_\Q \C$, deduit par extension des scalaires de $\Q$ a $E$ de
l'isomorphisme $\C^{\bar{S}} \xrightarrow{\sim} k_0 \otimes_\Q \C$. Ceci fournit pour $D(\chi)$, bien
defini mod $E^*$, un representant dans $(\Q \otimes_\Q \C)^* = \C^* \subset (E \otimes_\Q \C)^*$,
a savoir le determinant de l'inverse de la matrice $(\sigma_\alpha)$, pour $\sigma \in \bar{S}$ et
$\alpha$ parcourant une base de $k_0$ sur $\Q$. L'isomorphisme $\C^{\bar{S}} \xrightarrow{\sim} k_0 \otimes_\Q \C$
transforme la forme quadratique $\sum x_\sigma^2$ en la forme $\Tr(xy)$. Ceci permet
d'identifier $(\det(\sigma_\alpha))^2$ au discriminant de $k_0$:
\[
D(\chi) \sim \text{racine carree du discriminant de } k_0.
\]
\end{remark}

\subsection*{8.18}
\addcontentsline{toc}{subsection}{8.18}

Notons $p''(\chi; \sigma, \tau)$ l'image de $p'(\chi; \sigma, \tau)$ dans $\C^*/\Q^*$. Elle ne depend que
de $\chi$, $\sigma$ et $\tau$. Si un homomorphisme algebrique $\eta: k^* \to E^*$ verifie (8.2.1),
une de ses puissances est la partie algebrique d'un caractere de Hecke: $\eta^M = \chi_{\text{alg}}$.
De plus, si $\chi_{\text{alg}}' = \chi_{\text{alg}}''$, $\chi'$ et $\chi''$ ne different que par un caractere d'ordre
fini et $\chi'^M = \chi''^M$ pour $M$ convenable. On deduit de (8.1.1) que
$p''(\chi^M; \sigma, \tau) = p''(\chi; \sigma, \tau)^M$, et ceci permet de poser sans ambiguite
\[
p''(\eta; \sigma, \tau) = p''(\chi; \sigma, \tau)^{1/M}, \quad \text{pour } \eta^M = \chi_{\text{alg}}.
\]
Ces periodes obeissent au formalisme suivant:
\begin{align}
p''(\eta' \eta''; \sigma, \tau) &= p''(\eta'; \sigma, \tau) p''(\eta''; \sigma, \tau). \\
\intertext{(8.18.2) $p''(\chi; \sigma, \tau)$ ne change pas quand on remplace $E$ par une extension
$E'$ de $E$, et $\tau$ par un de ses prolongements a $E'$.}
\intertext{(8.18.3) $p''(\chi; \sigma, \tau)$ ne change pas quand on remplace $k$ par une extension
$k'$ de $k$, $\sigma$ par un de ses prolongements a $k'$, et $\chi$ par $\chi \circ N_{k'/k}$.}
\intertext{(8.18.4) Si $\alpha$ est un automorphisme de $k$, et $\beta$ un automorphisme de $E$, on a}
p''(\eta; \sigma, \tau) &= p''(\beta \eta \alpha^{-1}; \sigma \alpha^{-1}, \tau \beta^*). \\
\intertext{(8.18.5) Le complexe conjugue de $p''(\eta; \sigma, \tau)$ est $p''(\eta; c\sigma, c\tau)$.}
\intertext{(8.18.6) Pour $k = F = \Q$, $p''(\text{Id}; \text{Id}, \text{Id}) = 2\pi i$.}
\end{align}
Les formules (8.15.1) a (8.15.3) resultent de (8.1.1) a (8.1.3), (8.18.4) et (8.18.5)
se voient par transport de structure, et (8.18.6) resulte de (8.5).

\subsection*{8.19}
\addcontentsline{toc}{subsection}{8.19}

Soit $\eta: k^* \to E^*$ un homomorphisme verifiant (8.2.1). On suppose aussi
que $n(\eta; \sigma, \tau)$ ne vaut jamais $w/2$, ce qui permet de definir $(k \otimes_\Q E)_+$
comme en 8.15. Definissons $\eta^*: E^* \to k^*$ par
\[
\eta^*(y) = \det_{k \otimes E}(\cdot y, (k \otimes_\Q E)_+)^{-1}.
\]
Cet homomorphisme verifie encore (8.2.1), et
\[
n(\eta^*; \sigma, \tau) =
\begin{cases}
1 & \text{si } n(\eta; \sigma, \tau) < w/2,\\
0 & \text{si } n(\eta; \sigma, \tau) > w/2.
\end{cases}
\]
Si $\eta^*$ est la partie algebrique d'un caractere de Hecke $\chi^*$, $M(\chi^*)$ est le $H^1$
d'une variete abelienne sur $E$, a multiplication complexe par $k$, dont l'algebre de
Lie, comme $k \otimes_\Q E$-module, est isomorphe a $(k \otimes_\Q E)^{\eta^*}$.

\begin{proposition}\label{prop:8.20}
Avec les hypotheses et notations de 8.19, prenons pour $E$ un sous-corps de $\C$,
et notons $1$ l'inclusion identique de $E$ dans $\C$. On a
\[
p''(\eta; \sigma, 1) = \prod_{n(\sigma,1) < w/2} p''(\eta^*; 1, \sigma)^{n(\sigma,1) - w/2}.
\]
\end{proposition}

Si $\xi: E \hookrightarrow E' \subset \C$ est une extension finie de $E$, le membre de gauche ne
change pas quand on remplace $\eta$ par $\xi \circ \eta$ (8.18.2). On a $(\xi \circ \eta)^* = \eta^* \circ
N_{E'/E}$ et, par (8.18.3), le membre de droite ne change pas non plus.

Si $\xi: k \to k'$ est une extension finie de degre $d$ de $k$, le membre de gauche est
eleve a la puissance $d$ quand on remplace $k$ par $k'$, et $\eta$ par $\eta \circ N_{k'/k}$: un
plongement complexe de $k$ est induit par $d$ plongements de $k'$, et on applique
(8.18.3). De meme pour le membre de droite, par (8.18.2) et l'egalite
$(\eta \circ N_{k'/k})^* = d \cdot \eta^*$.

Ces compatibilites nous ramenent a supposer que $E$ est galoisien et que $k$ est
isomorphe a $E$. Pour chaque isomorphisme $\omega$ de $k$ dans $E$, posons
$n(\omega) = n(\omega; 1 \circ \omega, 1)$. Notant additivement le groupe des homomorphismes
de $k^*$ dans $E^*$, on a $\eta = \sum_\omega n(\omega) \omega$. Puisque
$n(\eta; 1 \circ \omega, 1) = n(\eta \circ \omega^{-1}; 1, 1)$ (8.15.4), on a
\begin{equation}\label{eq:8.20.1}
\prod_{n(\sigma,1) < w/2} p''(\eta; \sigma, 1) = \prod_{n(\omega) < w/2} p''(\eta \circ \omega^{-1}; 1, 1)
= p''\left(\sum_{n(\omega) < w/2} \eta \circ \omega^{-1}; 1, 1\right).
\end{equation}

Par ailleurs,
\[
\sum_{n(\omega) < w/2} \eta \circ \omega^{-1} = \sum_{\omega} \sum_{\substack{n(\omega_2) \\ n(\omega_2) < w/2}} n(\omega_2) \omega_2 \circ \omega^{-1}
= \sum_{\omega_2} n(\omega_2) \sum_{\substack{\omega \\ n(\omega) < w/2}} \omega_2 \circ \omega^{-1}
\]
\[
= \sum_{\omega_2} n(\omega_2) \cdot \omega_2 \circ \eta^*.
\]
Ceci permet de continuer (8.20.1) par
\[
= \prod_{\omega_2} p''(\omega_2 \circ \eta^*; 1, 1)^{n(\omega_2)} = \prod_{\omega_2} p''(\eta^*; 1, \omega_2)^{n(\omega_2)}
\]
(nouvelle application de 8.15.4), ce qui prouve 8.20.

\subsection*{8.21}
\addcontentsline{toc}{subsection}{8.21}

Combinant 8.16 et 8.20, on trouve pour la composante d'indice $1$ de
$c^+ R_{k/\Q} M(\chi) \in (E \otimes_\Q \C)^*/E^* = \C^{\Hom(E,\C)}/E^*$, l'expression suivante,
mod $\Q^*$:
\[
c^+ R_{k/\Q} M(\chi) \sim \prod_\sigma p''(\chi_{\text{alg}}; 1, \sigma)^{n(\sigma,1)}.
\]
Si $2w$ (i.e.~$R_{k/\Q} M(\chi)$) est critique, la Conjecture 2.8 affirme donc que
\[
L(1 \circ \chi, 0) \sim \prod_\sigma p''(\chi_{\text{alg}}; 1, \sigma)^{n(\sigma,1)} \pmod{\Q^*}.
\]
Pour $E$ assez grand, $\chi_{\text{alg}}$ est la partie algebrique d'un caractere de Hecke $\chi^*$,
et les periodes s'interpretent comme periodes d'integrales abeliennes (8.19), (8.3).
L'enonce obtenu est celui que Shimura a prouve pour $k$ de type CM, ou abelien sur
un corps de type CM (avec une restriction sur le poids).

\begin{remark}
Si $\eta: k^* \to E^*$ verifie (8.2.1), il existe des sous-corps $k'$ de $k$ et
$\xi: E' \to E$ de $E$, soit de type CM, soit egaux a $\Q$, et une factorisation
$\eta = \eta' \circ N_{k'/k}$. On a alors $p''(\eta; \sigma, \tau) = p''(\eta'; \sigma|_{k'}, \tau|_{E'})$.

Si maintenant $k$ et $E$ sont de type CM (ou $\Q$), et qu'on note encore $c$ leur
conjugaison complexe, on a $c\sigma = \sigma c$, $c\tau = \tau c$, et $\eta^c = c\eta$, d'ou
$p''(\eta; c\sigma, \tau) = p''(\eta^c; \sigma, \tau) = p''(\eta; \sigma, \tau)$ et $p''(\eta; \sigma, c\tau) =
p''(\eta; \sigma, \tau)$, d'ou $p''(\eta; \sigma, \tau) = p''(\eta; c\sigma, c\tau) =
\overline{p''(\eta; \sigma, \tau)}$: les periodes, a priori dans $\C^*/\Q^*$, sont reelles, i.e.,
dans $\R^*/(\Q^* \cap \R^*)$.
\end{remark}

\begin{remark}
Soient $G_1$ le groupe de Galois de la reunion des extensions de type CM de $\Q$
dans $\overline{\Q} \subset \C$, et $c \in G_1$ la conjugaison complexe. C'est un element central
de $G$. Si $\varphi$ est une fonction localement constante a valeurs entieres sur $G$, nous
poserons $\varphi^*(x) = \varphi(xc)$. Supposons que $\varphi + \varphi^*$ est constante. Soient $G_1$ un
quotient fini de $G$ tel que $\varphi$ se factorise par une fonction $\varphi_1$ sur $G_1$, et $k$ le corps
correspondant. L'hypothese faite signifie que l'endomorphisme $\sum \varphi_1(\sigma) \sigma$
de $k^*$ verifie (8.2.1). La periode $p''(\sum \varphi_1(\sigma) \sigma; 1, 1)$ ne depend pas du choix
de $G_1$; on pose $P(\varphi) = p''(\sum \varphi_1(\sigma) \sigma; 1, 1)$.

La fonctionnelle $P$ est un homomorphisme dans $\C^*/\Q^*$ du groupe des
fonctions localement constantes a valeurs entieres sur $G$ qui verifient
$\varphi + \varphi^* =$ constante.
\end{remark}


\section*{Appendix by N. Koblitz and A. Ogus}
\addcontentsline{toc}{section}{Appendix by N. Koblitz and A. Ogus}

Algebraicity of some products of values of the $\Gamma$ function.

Let $A_N = N^{-1}\Z/\Z - \{0\}$, and let $U_N = (\Z/N\Z)^*$ operate on $A_N$
in the obvious way. If $f: A_N \to \C$, define $\langle f \rangle: U_N \to \C$ by
$\langle f \rangle(u) = \sum_{a \in A_N} \langle a \rangle f(ua)$, where $\langle a \rangle$ denote the
representative of $a$ between 0 and 1.

We first compute
\[
H_0 = \{f: A_N \to \Q \mid \langle f \rangle \text{ est constante}\}.
\]

\begin{example}
If $n$ is a divisor of $N$ and $a \in A_N$, consider the set
$S_m = \{a + k/n: k = 0, 1, \ldots, n-1\} \cup \{-na\} \subset A_N$. It is easy to see
that its characteristic function is in $H_0$. Note that $S_1 = \{a, -a\}$.
\end{example}

\begin{proposition}
$H_0$ is generated by $\{S_m^a: n \mid N \text{ ou } n \text{ est premier et } na \neq 0\}$.
\end{proposition}

\begin{proof}
The orbits of $A_N$ under the action of $U_N$ correspond to the divisors $d$ of
$N$ with $1 < d \leq N$: each orbit can be written uniquely in the form
$U_N(1/d)$. The stabilizer subgroup $I_d$ of $1/d$ is
$\{u \in U_N: u \equiv 1 \mod d\}$ and the orbit of $1/d$ is canonically
isomorphic to $U_d$. Thus, a function $f$ on $A_N$ is determined by the collection
of functions $f_d: U_d \to \Q$ defined by $f_d(v) = f(v/d)$.

To prove the proposition, we first complexify, so as to be able to work with
characters of $U_N$.

\begin{lemma}
If $f: A_N \to \C$ and if $\chi$ is a character of $U_N$, then the inner product of
$\langle f \rangle$ and $\chi$ is given by:
\[
(\langle f \rangle \mid \chi) = -\sum_d L(0, \chi_d) |I_d| (f_d \mid \chi_d),
\]
where the sum is taken over those divisors $d$ of $N$ such that $\chi$ is pulled
back from a character $\chi_d$ of $U_d$.
\end{lemma}

\begin{proof}
We have
\[
(\langle f \rangle \mid \chi) = \sum_{u \in U_N} \langle f \rangle(u) \bar{\chi}(u)
= \sum_{a \in A_N} \sum_{u \in U_N} \langle a \rangle f(ua) \bar{\chi}(u)
\]
\[
= \sum_{d \mid N} \sum_{u \in U_d} \sum_{v \in U_N} \langle v/d \rangle f_d(udv) \bar{\chi}(v),
\]
where $u_d \in U_d$ is the image of $u$.

Now if we choose a set $U_d' \subset U_N$ of coset representatives for
$U_N \to U_d$, we can write:
\[
\sum_{u' \in U_d'} \sum_{v \in U_N} f_d(u_d' v) \bar{\chi}(v) \langle v \rangle
= \sum_{u' \in U_d'} \sum_{u'' \in U_d'} \sum_{w \in I_d} f_d(u_d' u'' w) \bar{\chi}(u'' w)
\]
\[
= \sum_{u'} \sum_{u''} f_d(u_d' u'') \bar{\chi}_d(u') \bar{\chi}(w) |I_d|
\]
Of course, this sum is zero unless $\chi$ is trivial on $I_d$, i.e., unless $\chi$ is pulled
back from a character $\chi_d$ of $U_d$, in which case it is
$\chi_d(u') |I_d| (f_d \mid \chi_d)$.

If we substitute this into the above expression for $(\langle f \rangle \mid \chi)$ and
note that $\langle v/d \rangle = |v|/d$, where $|v|$ is the smallest positive representative
of $v \in U_d$, we find:
\[
(\langle f \rangle \mid \chi) = -\sum_d L(0, \chi_d) |I_d| (f_d \mid \chi_d)
= -\sum_d L(0, \chi_d) |I_d| (f_d \mid \chi_d) \quad \text{as claimed.} \qedhere
\]
\end{proof}

To prove the proposition, we let $d_0$ be the largest divisor of $N$ such that
$f_{d_0} \not\equiv 0$; it suffices to prove that any $f$ in $H_0$ can be written as $g + f'$,
with $g$ a linear combination of the $S_m^a$'s and $f_{d'}' = 0$ for $d' > d_0$. Let
$d = d_0$; we shall show below that $f_d$ is a linear combination of functions which
factor through $U_d \to U_d/\{\pm 1\}$, and functions $h_p$ which factor through
$U_d \to U_{d/p}$, for some prime divisor $p$ of $d$, $p \neq d'$. Any function
of the first type is invariant under $\pm 1$, and hence the corresponding function on
the orbit $U_N(1/d)$ is a linear combination of $S_m^a$'s. A function of the second
type is invariant under the kernel $K$ of $U_d \to U_{d/p}$, which has order $p$ if
$p$ divides $d/p$ and order $p-1$ otherwise. In the first case,
$K = \{1 + kd/p: k = 0, \ldots, p-1\}$, and in
the second, it is this same set with one element deleted, namely, the value of
$1 + kd/p$ which is divisible by $p$. Thus, the set $S_{M}, M = (K \cdot w)/d$, with the
addition of one or two elements in the orbit $U_N(p/d)$. Hence it is clear that the
corresponding function on the orbit $U_N(1/d)$ can be written as a linear
combination of $S_{M'}^a$'s and functions supported on the orbit of $p/d$, and we
have obtained our desired decomposition $f = g + f'$.

It remains for us to prove the claim about $f_d$. If $d$ is twice an odd number the map
$U_d \to U_d/\{\pm 1\}$ is an isomorphism, and the claim is trivial. In the other
cases, choose a primitive odd character $\chi_d$ of $U_d$, and let $\chi$ be the pull-back
of $\chi_d$ to $U_N$. Since $\chi$ is nontrivial and $\langle f \rangle$ is constant, $(\langle f \rangle \mid \chi) = 0$.
Let us look at the expression for this inner product provided by the lemma; the
only nonzero terms come from those $f_e$ such that $f_e \not\equiv 0$ and $I_e \subseteq \Ker(\chi)$.
But then $\Ker(\chi)$ contains $I_e$ and $I_d$, hence also $I_{d/e}$, where $m = \gcd(d,e)$.
Since $\chi_d$ is primitive, $m = d$ and $d$ divides $e$, and since $f_e = 0$ for $e > d$,
$0 = (\langle f \rangle \mid \chi) = -L(0, \chi_d) |I_d| (f_d \mid \chi_d)$. Since $\chi_d$ is odd and primi-
tive, we conclude that $(f_d \mid \chi_d) = 0$, i.e., $f_d$ is orthogonal to every odd primitive
character of $U_d$. It follows that $f_d$ can be written as a linear combination of
characters which are even (factor through $U_d/\{\pm 1\}$) or imprimitive (factor
through some $U_{d/p}$). When $d = p$, we should also remark that the constant
functions on $U_p$ are also invariant under $\pm 1$, and hence are already covered
by the first case.
\end{proof}

\begin{theorem}
Suppose that $f: A_N \to \Q$ is such that $\langle f \rangle$ is constant:
$\sum_a \langle a \rangle f(ua) = k$ for all $u \in U_N$. Then
\[
P(f) =_{\text{def}} \prod_{a \in A_N} \Gamma(\langle a \rangle)^{f(a)}
\]
is an algebraic number.
\end{theorem}

\begin{proof}
The map $\langle \cdot \rangle: H_0 \to \Q$ sending any $f$ to $\sum_a \langle a \rangle f(a)$ is evidently linear,
and hence $P(f+g) = P(f) P(g)$, for $f$ and $g$ in $H_0$. By the proposition, any
$f$ in $H_0$ is a $\Q$-linear combination of the $S_m^a$'s; hence $nf$ is a sum of
e's for some $n$. Since $P(nf)^n = P(nf)$, we are reduced to checking the theorem
for the $S_m^a$'s described in the proposition. Noting that $\langle S_m^a \rangle = 1$ and that
$\langle S_p^a \rangle = (p+1)/2$ if $pa \neq 0$, one has only to appeal to the identities:
\[
\prod_{k=1}^{p-1} \Gamma\left(\frac{k}{p}\right) = p \cdot (2\pi)^{(p-1)/2} p^{-1/2},
\quad
\prod_{k=0}^{p-1} \Gamma\left(x + \frac{k}{p}\right) = (2\pi)^{(p-1)/2} p^{1/2-px} \Gamma(px). \qedhere
\]
\end{proof}

\begin{remark}
Kubert has recently obtained a much more precise result expressed in a
different terminology from which it follows that, if $H_0$ is $\Z$-valued, then $2f$
is a $\Z$-linear combination of the $S_m^a$'s (to appear).
\end{remark}


\begin{thebibliography}{99}
\bibitem{1}
M. V. Borovoi, Sur l'action du groupe de Galois sur les classes de cohomologie rationnelles
de type $(p,p)$ de varietes abeliennes, Mat. Sb. 94 (1974), 649--652. (Russian)

\bibitem{2}
J. Coates and S. Lichtenbaum, On $l$-adic zeta functions, Ann. of Math. (2) 98 (1973), 498--550.

\bibitem{3}
P. Deligne, Applications de la formule des traces aux sommes trigonometriques, dans SGA
4$^{1/2}$, 168--232.

\bibitem{4}
P. Deligne, Les constantes des equations fonctionnelles des fonctions $L$, Modular Functions
of One Variable. II, Lecture Notes in Math., vol. 349, Springer-Verlag, New York, 1973,
pp. 501--595.

\bibitem{5}
P. Deligne, Formes modulaires et representations $l$-adiques, Seminaire Bourbaki 355
(Fevrier 1969), Lecture Notes in Math., vol. 179, Springer-Verlag, New York, 139--172.

\bibitem{6}
P. Deligne, Theorie de Hodge. III, Inst. Hautes Etudes Sci. Publ. Math. 44 (1974), 5--77.

\bibitem{7}
B. Gross, On the periods of abelian integrals and a formula of Chowla and Selberg, Invent.
Math. 45 (1978), 193--211.

\bibitem{8}
Ju. I. Manin, Correspondences, motifs and monoidal transformations.

\bibitem{9}
Ju. I. Manin, Points paraboliques et fonction zeta des courbes modulaires, Izv. 36 (1972), 19--66. (Russian)

\bibitem{10}
I. I. Piateckii-Shapiro, Relations entre les conjectures de Tate et de Hodge pour les
varietes abeliennes, Mat. Sb. 85(4) (1971), 610--620 = Math. USSR-Sb. 14 (1971), 615--625.

\bibitem{11}
J. P. Serre, Abelian $l$-adic representations and elliptic curves (McGill), Benjamin, New York, 1968.

\bibitem{12}
J. P. Serre, Facteurs locaux des fonctions zeta des varietes algebriques (definitions et
conjectures), Sem. Delange-Pisot-Poitou 1969/70, expose 19.

\bibitem{13}
G. Shimura, On some arithmetic properties of modular forms of one and several variables,
Ann. of Math. (2) 102 (1975), 491--515.

\bibitem{14}
C. L. Siegel, Berechnung von Zetafunktionen an ganzzahligen Stellen, Nachr. Akad. Wiss.
Gottingen Math. Phys. Kl. II 10 (1969), 87--102.

\bibitem{15}
J. Tate, On the conjecture of Birch and Swinnerton-Dyer and a geometric analogue, Seminaire
Bourbaki 306 (1965/66), Benjamin, New York; reproduit dans Dix exposes sur la theorie des
schemas, North-Holland, Amsterdam, 1968.

\bibitem{16}
A. Weil, Jacobi sums as ``Grossencharaktere'', Trans. Amer. Math. Soc. 73 (1952), 487--491.

\bibitem{17}
A. Weil, Sommes de Jacobi et caracteres de Hecke, Nachr. Akad. Wiss. Gottingen Math.
Phys. Kl. II 1 (1974), 1--14.

\bibitem{18}
D. Zagier, Modular forms whose Fourier coefficients involve zeta functions of quadratic fields,
Modular Forms of One Variable. VI, Lecture Notes in Math., vol. 627, Springer-Verlag, New
York, 1977, pp. 105--169.

\bibitem{19}
N. Saavedra, Categories tannakiennes, Lecture Notes in Math., vol. 265, Springer-Verlag,
New York, 1972.
\end{thebibliography}

\vspace{1em}
\noindent
SGA. Seminaire de geometrie algebrique du Bois-Marie. Theorie des topos et
cohomologie etale des schemas.

\noindent
SGA 4, par M. Artin, A. Grothendieck et J. L. Verdier, Lecture Notes in Math.,
vols. 269, 270, 305, Springer-Verlag, New York.

\noindent
SGA 4$^{1/2}$, par P. Deligne, Cohomologie etale, Lecture Notes in Math., vol. 569,
Springer-Verlag, New York, 1977.

\vspace{2em}
\begin{flushright}
\textsc{Institut des Hautes Etudes Scientifiques}\\
91440 Bures-sur-Yvette, France
\end{flushright}

\end{document}